% =============================================================================
% Section: The Forced Derivation (34-Step, Steps 0--33)
% =============================================================================

\section{The Forced Derivation}
\label{sec:derivation}

This section presents the core contribution: a 34-step derivation (Steps 0--33) showing
that the structure of witnessable reasoning is \emph{forced}, not chosen.
Starting from absolute nothingness ($\noth$), we show that A0$^*$
(\cref{ax:a0star}), derived from $\noth$ via N0--N4, uniquely determines a carrier, a self-delimiting syntax, an
executability layer, an \emph{ordered} ledger carrying the witness algebra
$\WitAlg$, a truth quotient with observable-open topology,
time, energy, gauge invariance, a Myhill--Nerode congruence, witness-path
geometry, Bellman-optimal dynamics governed by the Coupling Law, a self-hosting closure, a universal
execution law with consciousness and trit-field valuation, and a K\"{a}hler
geometric structure on the space of truths.
The ordered witness algebra $\WitAlg$ is the foundational algebraic object:
its noncommutativity is internal from the start, making the multiset
(commutative) sector a derived quotient rather than a primitive.
Each step takes the form of a theorem with an explicit
\emph{forcing argument}: a proof that the step is the unique
A0$^*$-compliant continuation.  The derivation DAG is depicted in
\cref{fig:derivation-dag}.

\bigskip\noindent
\textbf{Notation.}  We use $\Carrier$, $\Del$, $\Pit$, $\Tim$, $\Ledger$,
$\UTM$, $\noth$, $\opnoth$, and the witness algebra $\WitAlg$ as defined in
\cref{sec:axioms}.  New objects introduced in this derivation---$\FinSlice$,
$\SdMap$, $\ObsOpen$, $\Eraser$, $\FixPt$, $\MNcong$, $\OrdLedger$,
$\WitCompose$, $\dsep$, $\dtime$, $\WitCat$, $\ExactGate$, $\EnergyLedger$,
$\ClMin$, $\CanEnum$, $\SelfHost$---are defined at their point of first
appearance.

\bigskip\noindent
\textbf{Structure.}  The derivation is organized into nine phases:
\begin{itemize}[nosep]
  \item \textbf{Phase~A} (Step~0): The only law---A0$^*$ and the forced output gate.
  \item \textbf{Phase~B} (Steps~1--2): Carrier and syntax.
  \item \textbf{Phase~C} (Steps~3--4): Executability.
  \item \textbf{Phase~D} (Steps~5--7): Recording and truth.
  \item \textbf{Phase~E} (Steps~8--11): Time, energy, and limits.
  \item \textbf{Phase~F} (Steps~12--17): Invariance and control.
  \item \textbf{Phase~G} (Steps~18--21): Solver with coupling law, self-hosting, and closure.
  \item \textbf{Phase~H} (Steps~22--25): Universal execution.
  \item \textbf{Phase~I} (Steps~26--33): Geometric reconstruction.
\end{itemize}

% ===========================================================================
\subsection{Phase A: The Only Law (Step 0)}
\label{sec:derivation:law}

\begin{theorem}[Step 0 --- Nothingness $\noth$ and the Forced Output Gate]
\label{thm:nothingness}
\emph{Nothingness} is the state with no admissible distinctions:
$\noth : \Del = \varnothing,\; \Carrier = \varnothing.$
A0$^*$ forces a trichotomous output gate: any query $q$ must terminate in exactly one of
$\UNIQUE \;(+1)$, $\OMEGA \;(0)$, or $\UNSAT \;(-1)$.
\begin{proof}[Proof sketch]
If no witness procedure exists then no distinction is admissible (A0$^*$), so $\noth$ is the unique structureless state.  After all affordable tests on a query with candidate set $\mathcal{A}$, the surviving set $\mathcal{A}_S$ has $|\mathcal{A}_S| = 1$ ($\UNIQUE$), $|\mathcal{A}_S| > 1$ ($\OMEGA$), or $|\mathcal{A}_S| = 0$ ($\UNSAT$)---three cases that partition the non-negative integers by pigeonhole.
\forcing{The trichotomy is forced by pigeonhole on $|\mathcal{A}_S| \in \{0, 1, 2, \ldots\}$ collapsed into three regimes. No fourth case exists.}
See Appendix~\ref{app:derivation} for the complete proof.
\end{proof}
\end{theorem}

% trit-states.tex — Three terminal states: UNIQUE, OMEGA, UNSAT
% Requires opoch.sty (loaded by main document)

\begin{figure}[htbp]
\centering
\begin{tikzpicture}[
  query/.style={diamond, draw=opochblue, fill=opochblue!10, minimum width=22mm,
    minimum height=18mm, font=\small\bfseries, inner sep=2pt, line width=0.8pt},
  tbox/.style={rounded rectangle, minimum width=38mm, minimum height=18mm,
    font=\small, inner sep=4pt, line width=0.8pt, align=center},
  ubox/.style={tbox, draw=opochgreen, fill=opochgreen!8},
  obox/.style={tbox, draw=opochorange, fill=opochorange!8},
  rbox/.style={tbox, draw=opochred, fill=opochred!8},
  tarr/.style={-{Stealth[length=3mm]}, thick},
  tlbl/.style={font=\footnotesize\bfseries, fill=white, inner sep=2pt, rounded corners=1pt},
]

  % Central query node
  \node[query] (q) at (0, 0) {Query $q$};

  % === UNIQUE (+1) — top right ===
  \node[ubox] (uniq) at (5.5, 2.2)
    {\textbf{\UNIQUE{} $(+1)$}\\[2pt]
     \scriptsize answer $+$ witness $+$ receipt};
  \draw[tarr, opochgreen!80] (q.east) -- ++(0.8, 0) |- (uniq.west)
    node[tlbl, pos=0.25, above] {$+1$};

  % Green check
  \node[font=\large, opochgreen!70!black] at (3.2, 2.2) {\checkmark};

  % === OMEGA (0) — middle right ===
  \node[obox] (omega) at (5.5, -0.0)
    {\textbf{$\OMEGA$ $(0)$}\\[2pt]
     \scriptsize family $+$ separator $+$ gap};
  \draw[tarr, opochorange!80] (q.east) -- (omega.west)
    node[tlbl, pos=0.5, above] {$0$};

  % Question mark
  \node[font=\large, opochorange!80!black] at (3.2, 0.0) {?};

  % === REFUTED (-1) — bottom right ===
  \node[rbox] (ref) at (5.5, -2.2)
    {\textbf{\REFUTED{} $(-1)$}\\[2pt]
     \scriptsize reason $+$ violations};
  \draw[tarr, opochred!80] (q.east) -- ++(0.8, 0) |- (ref.west)
    node[tlbl, pos=0.25, below] {$-1$};

  % X mark
  \node[font=\large\bfseries, opochred!70] at (3.2, -2.2) {$\times$};

  % Annotations on the left
  \node[font=\scriptsize\itshape, opochgray, align=right, anchor=east]
    at (-1.8, 1.0) {Every query\\terminates in\\exactly one of\\three states};

  % Trit label
  \node[font=\scriptsize, opochgray, align=center] at (5.5, -3.6)
    {$\trit{+1, 0, -1}$ --- the \emph{trit} replaces\\
     probabilistic confidence scores};

\end{tikzpicture}
\caption{The three terminal states forced by the output gate (Step~0).
  Every query~$q$ resolves to exactly one trit value:
  $\UNIQUE$ ($+1$) with a verified witness,
  $\OMEGA$ ($0$) with the minimal separator test needed to resolve,
  or $\UNSAT$ ($-1$) with specific violations identified.
  No probabilistic hedging is permitted.}
\label{fig:trit-states}
\end{figure}


% ===========================================================================
\subsection{Phase B: Carrier and Syntax (Steps 1--2)}
\label{sec:derivation:carrier}

\begin{theorem}[Step 1 --- Carrier $\Carrier$ with Finite Slice $\FinSlice$]
\label{thm:carrier}
A0$^*$ forces the carrier $\Carrier = \{0,1\}^{<\infty}$.  For budget $B$, the \emph{finite slice} is
$\FinSlice = \{ s \in \{0,1\}^* \mid |s| \leq B \}$.
\emph{Dependencies:} A0$^*$, Step~0.
\begin{proof}[Proof sketch]
A0$^*$ requires finite witness procedures, hence finite strings over some finite alphabet $\Sigma$.  Any $|\Sigma| \geq 2$ is computably isomorphic to $\{0,1\}$; unary admits no internal distinctions.  Thus $\Carrier = \{0,1\}^{<\infty}$ is ontologically unique up to computable isomorphism.  $\FinSlice$ is forced because a witness of cost $B$ cannot access strings longer than $B$ bits.
\forcing{Departure from $\noth$ requires at least one distinction, hence one witness, hence one finite string.  The set of all finite strings is $\Carrier$.  Budget-bounding gives $\FinSlice$.  (A0$^*$ forcing.)}
See Appendix~\ref{app:derivation} for the complete proof.
\end{proof}
\end{theorem}

\begin{theorem}[Step 2 --- Self-Delimiting Syntax $\SdMap$]
\label{thm:sd}
In a closed universe (no external delimiter oracle), descriptions must be self-delimiting.  The canonical self-delimiting map is
$\SdMap(s) = 1^{|s|} \, 0 \, s$.
\emph{Dependencies:} A0$^*$, Step~1.
\begin{proof}[Proof sketch]
Without self-delimitation, a prefix $x \sqsubseteq y$ creates an ontologically irresolvable ambiguity: no finite procedure can separate ``description is $x$'' from ``description begins with $x$,'' violating A0$^*$.  $\SdMap(s) = 1^{|s|}0s$ resolves this: read bits until the first $0$ to obtain $|s|$, then read $|s|$ bits.  The encoding is prefix-free; other self-delimiting encodings are computably isomorphic.
\forcing{Any non-self-delimiting encoding creates an A0$^*$-violating ambiguity.  Self-delimitation is the minimal constraint that eliminates it.}
See Appendix~\ref{app:derivation} for the complete proof.
\end{proof}
\end{theorem}

% ===========================================================================
\subsection{Phase C: Executability (Steps 3--4)}
\label{sec:derivation:executability}

\begin{theorem}[Step 3 --- Executability Primitives]
\label{thm:executability}
A0$^*$ forces three executability primitives (the \emph{declared substrate}):
(a)~a total evaluator $\UTM$ halting on all inputs,
(b)~a total decoder $\mathrm{dec} \colon \Carrier \to \Carrier \times \Carrier$,
(c)~a cost functional $c \colon \Del \to \mathbb{N}_{>0}$ with $c(\tau) \geq 1$.
\emph{Dependencies:} A0$^*$, Steps~1--2.
\begin{proof}[Proof sketch]
A0$^*$ demands every finite witness procedure be admissible, so $\Del$ must be maximally closed under totality-preserving composition---which is exactly the class of total computable functions (Proposition; proof in Appendix~\ref{app:derivation}).  A universal evaluator $\UTM$ is forced by the $s$-$m$-$n$ and universal function theorems \citep{rogers1967}.  The decoder follows from self-delimiting syntax (Step~2); the cost functional from A0$^*$'s finiteness requirement.
\forcing{(i)~A0$^*$ forces maximal closure of $\Del$ under totality-preserving composition.
(ii)~This class equals the total computable functions.
(iii)~$\UTM$ is forced.
(iv)~The form of $\UTM$ is gauge (Step~12).
(v)~CT enters only as a naming convention for the forced class.}
See Appendix~\ref{app:derivation} for the complete proof.
\end{proof}
\end{theorem}

\begin{remark}[CT as Naming Convention]
\label{rem:ct-localization}
The mathematical content of Step~3 is a proposition of this paper:
maximal closure under totality-preserving composition on a finite
alphabet is exactly the class of total computable functions (proof in
Appendix~\ref{app:derivation}).  The key identification of ``finite
procedure'' (informal, from A0$^*$) with ``total computable function''
(formal) is the Church--Turing thesis.

The Church--Turing thesis enters only as the identification of an
informal concept (``finite procedure,'' from A0$^*$'s natural-language
phrasing) with a formal one (``total computable function'').  This is
analogous to how ``number'' is identified with ``element of
$\mathbb{N}$''---a naming convention, not an empirical hypothesis.

The model choice (Turing machine vs.\ $\lambda$-calculus vs.\ register
machine) is gauge (formalized at Step~12): the truth quotient is
invariant under switching between isomorphic computation models.  A
reader who replaces ``Turing-computable'' with a weaker class (e.g.,
primitive recursive functions) obtains a weaker but structurally
identical derivation.

Thus Step~3 is forced: A0$^*$'s maximal closure yields the total computable
functions (see Appendix~\ref{app:derivation} for the proof; general
background in \citealt{rogers1967}, Chapters~1--2).  The Church--Turing
thesis corroborates this conclusion; it does not enter the derivation
as a premise.
\end{remark}

\begin{remark}[Total Semantics]
\label{rem:total-semantics}
The decoder $\mathrm{dec}$ and evaluator $\UTM$ are total by
construction: every input produces a well-defined output.  In
particular, the outcome alphabet $A$ of any test $\tau$ includes
explicit $\mathrm{FAIL}$ and $\mathrm{TIMEOUT}$ outcomes.
``Undefined behavior'' is not a gap in the semantics but a
specific outcome in~$A$: if $\UTM$ exceeds the allotted budget,
it returns $\mathrm{TIMEOUT} \in A$, not $\bot$.  This ensures
that no test can create a semantic gap---every execution path
terminates in a recorded outcome, as required by A0$^*$.
\end{remark}

\begin{theorem}[Step 4 --- Endogenous Tests $\Del^*$ and Canonical Enumeration $\CanEnum$]
\label{thm:endogenous}
The test algebra is self-generating: $\Del^* = \mathrm{lfp}(\ClMin)(\Del_0)$.
The canonical enumeration $\CanEnum \colon \mathbb{N} \to \Del^*$ lists tests by $\SdMap$-length (ties broken lexicographically), requiring no scheduler channel.
\emph{Dependencies:} A0$^*$, Steps~1--3.
\begin{proof}[Proof sketch]
The closure rules (composition, branching, bounded iteration) are each individually forced by A0$^*$: omitting any one excludes an admissible finite procedure.  These rules are monotone on the powerset lattice, so by Knaster--Tarski the least fixed point exists.  The enumeration order is forced by $\SdMap$: any other ordering would require an external channel, violating closure.
\forcing{Any subset not closed under $\ClMin$ excludes a valid test---an A0$^*$-violating distinction.  The least fixed point is the minimal set avoiding this.}
See Appendix~\ref{app:derivation} for the complete proof.
\end{proof}
\end{theorem}

% ===========================================================================
\subsection{Phase D: Recording and Truth (Steps 5--7)}
\label{sec:derivation:recording}

\begin{theorem}[Step 5 --- Ordered Ledger $\OrdLedger$ and the Diamond Law as Theorem]
\label{thm:ledger}
Witnesses must persist; therefore an append-only ledger is forced.  Under A0$^*$, the
primitive ledger is \emph{ordered}:
\[
  \OrdLedger \;=\; (e_1 \WitCompose e_2 \WitCompose \cdots \WitCompose e_n),
\]
where $\WitCompose$ is ordered witness composition in $\WitAlg$ and each entry
$e_i = (\tau_i,\, a_i,\, c(\tau_i),\, \delta_i)$ records the test, outcome,
cost, and state-update.  Append-only persistence is still forced: deleting
$e_i$ destroys replay of every witness that causally depends on~$e_i$.

\medskip\noindent
\textbf{Witness commutator.}  For any two entries $e_i, e_j$ define the
\emph{witness commutator}
\[
  \WitCommutator{e_i}{e_j}
    \;=\; \mathrm{effect}(e_i \WitCompose e_j)
         - \mathrm{effect}(e_j \WitCompose e_i).
\]

\medskip\noindent
\textbf{Diamond Law (Theorem).}
When $\WitCommutator{e_i}{e_j} = 0$, the order of $e_i$ and $e_j$ does not
affect any downstream witness; hence those pairs quotient to a multiset.
The Diamond Law is therefore a \emph{theorem about the commuting sector}
of $\WitAlg$, not a global assumption:
\begin{quote}
For $e_i, e_j$ with $\WitCommutator{e_i}{e_j} = 0$: \;
$\Noise \circ \Query = \Query \circ \Noise$ on $\TQ{\OrdLedger}$.
\end{quote}
The multiset ledger $\Ledger = \{\!\!\{(\tau_i, a_i)\}\!\!\}$ of the
original formulation is recovered as the quotient of $\OrdLedger$ by the
commuting sector.
\emph{Dependencies:} A0$^*$, Steps~1--4.
\Cref{fig:diamond-law} illustrates the Diamond Law's commutativity of independent ledger entries.
\begin{proof}[Proof sketch]
If a witness's output could be erased, subsequent tests depending on it would
lack their input, making the distinction unwitnessable---violating A0$^*$.
Hence an append-only record is forced.  A0$^*$ supplies the ordered witness
algebra $\WitAlg$, so the primitive record inherits its ordering:
$\OrdLedger = (e_1 \WitCompose \cdots \WitCompose e_n)$.  Each entry records
the full 4-tuple $(\tau_i, a_i, c(\tau_i), \delta_i)$ because omitting any
component loses a witnessed datum.

For the Diamond Law: if $\WitCommutator{e_i}{e_j} = 0$ then swapping $e_i$
and $e_j$ in $\OrdLedger$ produces an identical downstream effect---every
future witness is replay-equivalent.  Order between such pairs is therefore
ontologically untestable and quotients away by the truth quotient (Step~6).
When $\WitCommutator{e_i}{e_j} \neq 0$, order matters and the
noncommutative structure is physical.
\forcing{Any mutable record allows A0$^*$ violations via witness destruction.
Ordered structure is forced because $\WitAlg$ is noncommutative.  The Diamond
Law is a theorem about the commuting sub-algebra, not a postulate.}
See Appendix~\ref{app:derivation} for the complete proof.
\end{proof}
\end{theorem}

\begin{theorem}[Step 6 --- Truth Quotient $\TQ{\Ledger}$]
\label{thm:truth-quotient}
Descriptions indistinguishable by all tests in the ledger are identified.  The resulting quotient $\TQ{\Ledger} = \FinSlice / {\equivL}$ is the truth structure.
\emph{Dependencies:} A0$^*$, Steps~1--5.
\begin{proof}[Proof sketch]
Define ledger equivalence: $x \equivL y$ iff no recorded test distinguishes them.  This is an equivalence relation by inspection.  Maintaining distinctions between elements that no test can separate would assert a witnessless distinction, violating A0$^*$.  For truth class $p$, the fiber $\fiber{p} = \{x \in \FinSlice \mid [x]_{\equivL} = p\}$ captures the gauge redundancy.
\forcing{Maintaining distinctions when no test can separate them violates A0$^*$.  Quotienting is forced.}
See Appendix~\ref{app:derivation} for the complete proof.
\end{proof}
\end{theorem}

\begin{remark}[Fiber Intuition]
\label{rem:fiber-intuition}
The fiber $\fiber{p}$ is the set of descriptions that ``mean the same thing''---all
bit-level encodings that no test can tell apart.  Quotienting by fibers is
the formal version of ignoring irrelevant notation.
\end{remark}

\begin{theorem}[Step 7 --- Indistinguishability as Primary Algebra]
\label{thm:indistinguishability}
The truth quotient carries three structures forced by A0$^*$: \emph{observable opens} $\ObsOpen$, an \emph{eraser algebra} $\Eraser$, and a \emph{self-generating fixed point} $\mathrm{sim}^*$.
\emph{Dependencies:} A0$^*$, Steps~4--6.
The observable-open topology is depicted in \cref{fig:observable-opens}.
\begin{proof}[Proof sketch]
(a)~Observable opens $\ObsOpen = \{O_{\tau,a}\}$ form a topology on $\TQ{\Ledger}$: each $O_{\tau,a}$ is the set of truth classes producing outcome $a$ under test $\tau$, closed under finite intersection and arbitrary union.  (b)~The eraser algebra $\Eraser = \{e_\tau\}$ consists of idempotent, monotone, commuting surjections that merge indistinguishable classes.  (c)~$\mathrm{sim}^* = \mathrm{lfp}(\Phi)$ captures self-referential closure: truths that survive are exactly those distinguishable by tests drawn from the truths themselves.  Existence via Knaster--Tarski on the finite lattice $2^{\TQ{\Ledger}}$.
\forcing{$\ObsOpen$ is forced because A0$^*$ equates meaningful with testable distinction.  $\Eraser$ is forced as the canonical idempotent operation on truth quotients.  The fixed point is forced because the system must be consistent with its own outputs acting as inputs.}
See Appendix~\ref{app:derivation} for the complete proof.
\end{proof}
\end{theorem}

% ===========================================================================
\subsection{Phase E: Time, Energy, and Limits (Steps 8--11)}
\label{sec:derivation:time}

\begin{theorem}[Step 8 --- Observer Collapse and Entropic Time]
\label{thm:time}
Time is the irreversible formation of records.  Each new record collapses the fiber:
$\Del\Tim = \log(|\fiber{}|_{\mathrm{pre}} / |\fiber{}|_{\mathrm{post}}) \geq 0$.
\emph{Dependencies:} Steps~5--7.
\Cref{fig:fiber-shrinkage} depicts the monotone fiber shrinkage under successive tests.
\begin{proof}[Proof sketch]
Each new entry $(\tau, a)$ can only refine the truth quotient (never coarsen) because the ledger is append-only (Step~5) and all tests are total (Step~3).  Hence $|\fiber{}|_{\mathrm{post}} \leq |\fiber{}|_{\mathrm{pre}}$ and $\Del\Tim \geq 0$.  This is a Second Law: time \emph{is} the irreversibility of witnessing.
\forcing{Reversing a ledger entry would destroy a witness (Step~5), violating A0$^*$.  Time \emph{is} the irreversibility of record formation.}
See Appendix~\ref{app:derivation} for the complete proof.
\end{proof}
\end{theorem}

\begin{theorem}[Step 9 --- Entropy and Budget]
\label{thm:entropy}
The entropy of the truth quotient is $S_p(\Ledger) = \log |\fiber{p}|$, and the total budget bounds affordable tests: $\sum_i c(\tau_i) \leq B$.
\emph{Dependencies:} Steps~5, 8.
\begin{proof}[Proof sketch]
Entropy decreases with each informative test (Step~8).  A0$^*$'s finiteness requirement forces a budget $B$.  The feasible test set $\Del(\Tim)$ is monotonically non-increasing---an expanding set would require infinite resources at finite time, violating A0$^*$.
\forcing{An expanding test set would require infinite resources at finite time---violating A0$^*$.  The budget \emph{is} the entropy available for allocation.}
See Appendix~\ref{app:derivation} for the complete proof.
\end{proof}
\end{theorem}

\begin{theorem}[Step 10 --- Energy Ledger $\EnergyLedger$]
\label{thm:energy}
Each test incurs energy cost $\Del E = c(\tau)$.  Efficiency: $\Efficiency = \Del\Tim / c(\tau)$.
\emph{Dependencies:} Steps~8--9.
\begin{proof}[Proof sketch]
By Landauer's principle, erasing information requires energy dissipation; A0$^*$'s append-only ledger makes record formation irreversible.  The energy ledger $\EnergyLedger = [(c(\tau_i), \Del\Tim_i), \ldots]$ is forced because A0$^*$ demands finite recording of each distinction's cost; the efficiency ratio is the canonical dimensionless measure.
\forcing{The energy ledger is forced because A0$^*$ demands finite cost recording.  The efficiency ratio is canonical from the two quantities the framework produces.}
See Appendix~\ref{app:derivation} for the complete proof.
\end{proof}
\end{theorem}

\begin{theorem}[Step 11 --- Operational Nothingness $\opnoth$]
\label{thm:opnoth}
When the budget is exhausted, the feasible test set becomes empty ($\Del(\Tim_{\max}) = \varnothing$): operational nothingness.
\emph{Dependencies:} Steps~8--10.
\begin{proof}[Proof sketch]
As $B_{\mathrm{spent}} \to B$, the feasible test set shrinks to the constant function $\tau_0(x) = c$, which distinguishes nothing.  $\opnoth$ is the unique terminal state under A0$^*$.  Unlike $\noth$, it occurs \emph{within} a structured universe---the horizon of resolvability.
\forcing{Finite budget $\Rightarrow$ finite test sequence $\Rightarrow$ terminal state.  $\opnoth$ is unique under A0$^*$.}
See Appendix~\ref{app:derivation} for the complete proof.
\end{proof}
\end{theorem}

% ===========================================================================
\subsection{Phase F: Invariance and Control (Steps 12--17)}
\label{sec:derivation:invariance}

\begin{theorem}[Step 12 --- Gauge: No Label Privilege]
\label{thm:gauge}
The set of untestable relabelings forms a group $\gauge \cong \prod_p \mathrm{Sym}(\fiber{p})$.
\emph{Dependencies:} Steps~6, 9.
\Cref{fig:gauge-filter} illustrates how gauge transformations permute within fibers without affecting the truth quotient.
\begin{proof}[Proof sketch]
A relabeling $g \colon \FinSlice \to \FinSlice$ permuting within fibers ($g(\fiber{p}) = \fiber{p}$) is untestable under A0$^*$.  The collection $\gauge$ is closed under identity (doing nothing is undetectable), inverse ($g^{-1} \circ g = \mathrm{id}$), and composition (C1).  Since each $g$ is a bijection on all of $\FinSlice$, composition is always defined, giving a group (the direct product of symmetric groups on fibers).
\forcing{Each group axiom is individually forced: omitting identity, inverse, or composition from $\gauge$ would render an untestable transformation detectable, contradicting the truth quotient.}
See Appendix~\ref{app:derivation} for the complete proof.
\end{proof}
\end{theorem}

\begin{remark}[Gauge Intuition]
\label{rem:gauge-intuition}
The gauge group $\gauge$ is the set of relabelings that no experiment can
detect---permuting descriptions within a fiber changes nothing observable.
Physical gauge invariance (electromagnetism, general covariance) is a
sector-specific instance of this universal principle.
\end{remark}

\begin{theorem}[Step 13 --- Decision Universes: Epistemic vs.\ Decision Regimes]
\label{thm:regimes}
The system operates in one of two regimes: \emph{epistemic} $\Epistemic$ (gathering information) or \emph{decision} $\Decision$ (committing to $\UNIQUE$ or $\OMEGA$).
\emph{Dependencies:} Steps~6, 8, 12.
\begin{proof}[Proof sketch]
The system must either execute another test (epistemic: fiber shrinks by evidence) or declare a terminal state (decision: irreversible commitment).  Conflating modes would allow asserting answers without evidence or gathering evidence without recording it---both A0$^*$ violations.  Transition $\Epistemic \to \Decision$ occurs when $|\mathcal{A}_S| \leq 1$ or the budget is exhausted.
\forcing{Conflating epistemic and decision modes permits unwitnessed assertions or destroyed witnesses.  Both violate A0$^*$.}
See Appendix~\ref{app:derivation} for the complete proof.
\end{proof}
\end{theorem}

\begin{theorem}[Step 14 --- Exactness Gate $\ExactGate$ and Split Law]
\label{thm:exactness}
Every query is processed \emph{certificate-first}: check the ledger for a witness ($\UNIQUE$) or counter-witness ($\UNSAT$) before searching for a separator.  Every transformation factors as merge $\circ$ relabel.
\emph{Dependencies:} Steps~5--6, 13.
\begin{proof}[Proof sketch]
Certificate-first ordering is forced by budget optimality: searching when a certificate exists wastes budget; declaring $\OMEGA$ when a certificate exists withholds an admissible witness.  For the Split Law: any total computable $\bar{f}$ on truth classes factors as $\bar{f} = \iota \circ q$, separating information loss from gauge relabeling.
\forcing{Certificate-first is forced by budget optimality.  The Split Law is forced because non-factorable transformations conflate information-destroying and information-preserving components.}
See Appendix~\ref{app:derivation} for the complete proof.
\end{proof}
\end{theorem}

\begin{theorem}[Step 15 --- Myhill--Nerode Congruence $\MNcong$]
\label{thm:myhill-nerode}
Two states $s, s'$ are Myhill--Nerode equivalent ($s \MNcong s'$) iff all futures agree:
$s \MNcong s' \iff \forall\, \sigma \in \Del^{*\,\omega}: \mathrm{outcome}(s, \sigma) = \mathrm{outcome}(s', \sigma)$.
$\MNcong$ is a right-congruence on the monoid $(\Del^*, \circ)$.
\emph{Dependencies:} Steps~4, 6, 12.
\Cref{fig:myhill-nerode} illustrates the Myhill--Nerode congruence classes and the resulting minimal automaton.
\begin{proof}[Proof sketch]
The state $s = (\TQ{\Ledger}, \Del(\Tim), B - B_{\mathrm{spent}})$ is futures-equivalent to $s'$ when no continuation distinguishes them.  Right-congruence: if $s \MNcong s'$ then $s \cdot \tau \MNcong s' \cdot \tau$.  The quotient $\Del^{*\,\omega}/{\MNcong}$ is finite (since $\FinSlice$ is finite), yielding a unique minimal automaton by the classical Myhill--Nerode theorem.
\forcing{Maintaining a distinction when all futures agree violates A0$^*$.  The Myhill--Nerode quotient is the coarsest A0$^*$-consistent partition, hence forced.}
See Appendix~\ref{app:derivation} for the complete proof.
\end{proof}
\end{theorem}

\begin{theorem}[Step 16 --- Witness-Path Geometry ($\dsep$, $\dtime$)]
\label{thm:separator-geometry}
The truth classes carry a witness-path geometry consisting of a symmetric
metric $\dsep$ and a directed quasi-metric $\dtime$, both induced by the
witness path category.
\emph{Dependencies:} Steps~4, 6, 15.
The separator geometry is depicted in \cref{fig:separator-geometry}.

\medskip\noindent
\textbf{Witness path category.}  Define $\WitCat$ with truth classes as
objects and replayable witness paths $\WitPath \colon p \to q$ as morphisms.
A morphism $\WitPath$ is a finite sequence of tests in $\OrdLedger$ that
transforms $p$ into $q$; composition is path concatenation.

\medskip\noindent
\textbf{Witness cost split.}  Every witness path $\WitPath$ decomposes as
\[
  K(\WitPath) \;=\; \Ksep(\WitPath) \;+\; \Kdiss(\WitPath),
\]
where $\Ksep(\WitPath)$ is the \emph{reversible separative cost}
(the cost of distinguishing endpoint truth classes, invariant under
output relabeling) and $\Kdiss(\WitPath)$ is the \emph{irreversible
disturbance cost} (the entropy produced by the path, directed and
non-negative by Step~8).

\medskip\noindent
\textbf{Symmetric metric.}  Define
\[
  \dsep(p, q) \;=\; \inf\bigl\{\Ksep(\WitPath) : \WitPath \colon p \to q
  \text{ in } \WitCat\bigr\}.
\]
This is a \emph{symmetric} metric: separating $p$ from $q$ and $q$ from $p$
differ only by output relabeling, which is gauge (Step~12).  The triangle
inequality is \emph{automatic} from path concatenation: for any
$p, q, r$, the concatenation $\WitPath_{p \to q} \cdot \WitPath_{q \to r}$
is a valid path $p \to r$, so
$\dsep(p, r) \leq \Ksep(\WitPath_{p \to q}) + \Ksep(\WitPath_{q \to r})$;
taking infima yields $\dsep(p, r) \leq \dsep(p, q) + \dsep(q, r)$.  No
transitively-separating property~(TS) is needed.

\medskip\noindent
\textbf{Directed quasi-metric.}  Define
\[
  \dtime(p, q) \;=\; \inf\bigl\{\Kdiss(\WitPath) : \WitPath \colon p \to q
  \text{ in } \WitCat\bigr\}.
\]
This is a directed quasi-metric ($\dtime(p,q) \neq \dtime(q,p)$ in general)
capturing time and thermodynamics: irreversible disturbance is directional.

\medskip\noindent
\textbf{Ordered witness geometry.}  The pair $(\dsep, \dtime)$ is the
forced geometry on truth classes.  Curvature $\SepCurv$ and hardness
$\Hardness$ are derived from $\dsep$ as before.
\begin{proof}[Proof sketch]
The witness path category $\WitCat$ is forced by A0$^*$: truth classes are
the objects of the theory (Step~6), and replayable witness paths are the
morphisms forced by the ordered ledger (Step~5).  The cost split
$K = \Ksep + \Kdiss$ is forced because A0$^*$'s witness algebra $\WitAlg$
distinguishes reversible and irreversible components: $\Ksep$ measures
the minimum cost of \emph{any} separating path (invariant under output
relabeling, hence symmetric), while $\Kdiss$ measures the entropy
produced (irreversible by Step~8, hence directed).

Symmetry of $\dsep$: swapping the roles of $p$ and $q$ in a separating
test amounts to relabeling outputs, which is gauge (Step~12);
hence $\dsep(p,q) = \dsep(q,p)$.  The triangle inequality follows from
path concatenation in $\WitCat$ without any additional property.

Non-symmetry of $\dtime$: the disturbance cost along $\WitPath: p \to q$
generally differs from the cost along any path $q \to p$ because entropy
production is path-dependent and directional (Step~8).
\forcing{A0$^*$ defines distinguishability via tests with costs in $\WitAlg$.
The witness path category is forced by the ordered ledger.  The cost split
into $\Ksep$ and $\Kdiss$ is forced by the reversible/irreversible
decomposition.  Symmetry of $\dsep$ is forced by gauge invariance; the
triangle inequality is automatic from path concatenation.  No external
structural commitment (TS) is needed.}
See Appendix~\ref{app:derivation} for the complete proof.
\end{proof}
\end{theorem}

\begin{theorem}[Step 17 --- $\Pit$-Consistent Control]
\label{thm:pi-consistent}
The truth map $\Pit$ satisfies the idempotency law:
$\Pit \circ \Noise = \Pit \circ \Noise \circ \Pit = \Pit$,
where $\Noise \in \gauge$.
\emph{Dependencies:} Steps~6, 12.
\begin{proof}[Proof sketch]
Gauge transformation $\Noise$ permutes within fibers: $\Noise(x) \in \fiber{[x]}$, hence $\Pit(\Noise(x)) = \Pit(x)$, giving $\Pit \circ \Noise = \Pit$.  Further: $\Pit \circ \Noise \circ \Pit = \Pit^2 = \Pit$ by idempotence.  If $\Pit$ were gauge-sensitive, it would distinguish gauge-equivalent descriptions, contradicting the truth quotient.
\forcing{Gauge-sensitivity of $\Pit$ would constitute a test separating equivalent descriptions---contradicting the truth quotient.  $\Pit$-consistency is the operational expression of gauge invariance.}
See Appendix~\ref{app:derivation} for the complete proof.
\end{proof}
\end{theorem}

\begin{remark}[$\Pit$-Consistency as Architectural Bridge]
\label{rem:pi-consistency-bridge}
$\Pit$-consistency means that all control decisions depend only on
the truth quotient, never on raw descriptions.  This property is the
bridge between gauge invariance (Step~12) and the Bellman solver
(Step~18): the solver's inputs are guaranteed gauge-free, ensuring that
test selection is invariant under relabeling.  Algebraically, the
idempotency law $\Pit \circ \Noise = \Pit$ has the same form as a
retraction in category theory, connecting to the eraser algebra
(Step~7): each eraser $e_\tau$ is itself a retraction, and $\Pit$ is
the universal such retraction that factors through all of them.
\end{remark}

\begin{definition}[$\PiCanon$-Canonicalization]
\label{def:pi-canon}
Let $\Ser \colon \FinSlice \to \{0,1\}^*$ be the serialization induced by
$\SdMap$, and let $\equivL$ be the ledger-induced equivalence
(Definition~\ref{def:prim-pi}).  The \emph{$\Pit$-canonicalization} of
$x \in \FinSlice$ is:
\[
  \PiCanon(x) \;=\; \min_{\mathrm{lex}}\bigl\{\,\Ser(y) \;\bigm|\; y \equivL x\,\bigr\}.
\]
$\PiCanon$ satisfies:
\begin{enumerate}[nosep]
  \item \emph{Totality}: $\PiCanon(x)$ is defined for every $x \in \FinSlice$
        (the equivalence class is non-empty and finite).
  \item \emph{Slack collapse}: $\PiCanon(x) = \PiCanon(y)$ iff $x \equivL y$.
  \item \emph{Idempotence}: $\PiCanon(\PiCanon(x)) = \PiCanon(x)$.
  \item \emph{Gauge absorption}: for every $g \in \gauge$,
        $\PiCanon(g \cdot x) = \PiCanon(x)$.
\end{enumerate}
\end{definition}

% ===========================================================================
\subsection{Phase G: Solver, Self-Hosting, and Closure (Steps 18--21)}
\label{sec:derivation:solver}

\begin{theorem}[Step 18 --- Deterministic Solver (Prior-Free Minimax and Coupling Law)]
\label{thm:bellman}
The optimal test selection policy is a prior-free minimax Bellman recursion,
upgraded by the \emph{Separator-Indistinguishability Coupling Law} to a fully
deterministic, question-relative selector.
\emph{Dependencies:} Steps~9, 14, 16.

\medskip\noindent
\textbf{Part (a): Forced minimax (not $L^2$).}
Before any prior is itself witnessed and recorded in~$\OrdLedger$, branch
labels over future outcomes are gauge: no unrecorded probability distribution
over outcomes is admissible by A0$^*$.  Hence any legal policy must be
invariant under:
\begin{enumerate}[nosep]
  \item relabeling of unwitnessed outcomes,
  \item positive rescaling of cost,
  \item replay-equivalent decomposition of paths in $\WitCat$,
  \item quotient-preserving recodings (gauge, Step~12).
\end{enumerate}
The only cost aggregator respecting all four symmetries without a witnessed
prior is the \emph{worst-case continuation cost}---i.e., minimax.  Any
$L^2$ (mean-squared) or expected-value aggregator requires a probability
measure over branches, which is itself a distinction that must be witnessed
before use; importing an unwitnessed measure violates A0$^*$.

The resulting value function satisfies
$\tau^* = \argmin_\tau [c(\tau) + \max_a V(\OrdLedger \WitCompose (\tau,a),
B - c(\tau))]$.
Well-defined because $\Del(\Tim)$ is finite with positive integer costs.
Termination by induction on $(|\TQ{\OrdLedger}|_{\mathcal{A}}, B)$ under
lexicographic order.

\medskip\noindent
\textbf{Part (a$'$): Bayesian selection after witnessed prior.}
Once a prior $\mu$ has itself been witnessed (i.e., $\mu$ appears as a
recorded datum in $\OrdLedger$ with witness status), Bayesian
expected-cost selection becomes admissible---the prior is no longer an
unwitnessed externality.  Minimax is therefore the \emph{bootstrap}
policy; Bayesian selection is the \emph{mature} policy enabled by
self-observation.

\medskip\noindent
\textbf{Part (b): Coupling Law.}
For query $q$ and surviving answer set $W$, define the \emph{answer quotient}
$\AnsQuot{q}{W_a}$, the \emph{refinement gain}
$\RefGain{q}{\tau;\, W} = \log|\AnsQuot{q}{W}| - \max_a \log|\AnsQuot{q}{W_a}|$,
and the \emph{question-relative efficiency}
$\EffQ{\tau}{W} = \RefGain{q}{\tau;\, W} \,/\, c(\tau)$.
The \emph{Coupling Law} selects the next test as:
\[
  \boxed{\;
    \OptTest{q}{W} \;=\;
    \PiCanon\text{-}\!\min\Bigl(
      \argmax_{\tau \in \Del_{\leq B}}\;
      \EffQ{\tau}{W}
    \Bigr)
  \;}
\]
where $\PiCanon$-$\min$ is the $\Pit$-canonical tie-break (lexicographically
minimal under $\CanEnum$ order after gauge reduction).

\medskip\noindent
\textbf{Part (c): Deterministic evolution.}
Since $\EffQ{\tau}{W}$ is a deterministic function of $(W, q)$ alone,
and $\PiCanon$-$\min$ eliminates all tie-breaking freedom, the entire
epistemic evolution $W_0 \to W_1 \to \cdots$ is fully determined by the
initial pair $(q, W_0)$.  The \emph{Omega distance}
$\OmegaDist{q}{W} = \log|\AnsQuot{q}{W}|$ measures how far the current
state is from $\UNIQUE$.

\begin{proof}[Proof sketch]
Part~(a): the four symmetries (outcome relabeling, cost rescaling,
replay-equivalent decomposition, gauge recoding) each individually exclude
any aggregator that weights branches unequally without a witnessed basis for
the weighting.  The intersection of all four invariance groups leaves exactly
the worst-case (minimax) aggregator.  In particular, $L^2$ aggregation
requires a probability measure to define the expectation, and no such measure
exists in $\OrdLedger$ before it is itself witnessed---so $L^2$ is
inadmissible at the prior-free stage.
Part~(a$'$): once $\mu$ is recorded, the symmetry under outcome relabeling
is broken by $\mu$'s weights, and Bayesian selection becomes the unique
$\mu$-compatible policy.
Part~(b): the efficiency $\EffQ{\tau}{W}$ is the canonical dimensionless ratio
of information gained to cost paid, both quantities already present in the
framework (Steps~8, 10).  The $\argmax$ selects the most cost-effective
test; $\PiCanon$-$\min$ is forced by gauge invariance (Step~12).
Part~(c): since both $\varepsilon$ and the tie-break are deterministic
functions of $(W, q)$, the evolution is deterministic.
\forcing{Minimax is forced by A0$^*$: prior-free branch labels are gauge,
and the only aggregator invariant under all four symmetries without a
witnessed prior is worst-case continuation cost.  Bayesian selection
appears only after a prior itself has witness status in the trace.
The efficiency functional is uniquely determined: any efficiency
functional satisfying scale-invariance and monotonicity is of the form
$f(g/c)$ for strictly monotone~$f$, and $\argmax_\tau f(g(\tau)/c(\tau))
= \argmax_\tau g(\tau)/c(\tau)$ (see Remark~\ref{rem:efficiency-uniqueness}).
$\PiCanon$ tie-breaking is forced by gauge
invariance.  Together, they eliminate the last hidden degree of freedom
(the scheduler of witnessing), upgrading the framework from consistent
model to fully deterministic self-computing system.}
See Appendix~\ref{app:derivation} for the complete proof.
\end{proof}
\end{theorem}

% bellman-coupling.tex — Bellman minimax with Coupling Law (Step 18)
% Requires opoch.sty (loaded by main document)

\begin{figure}[ht]
\centering
\begin{tikzpicture}[
  qnode/.style={circle, draw=opochblue, fill=opochblue!12, minimum size=10mm,
    font=\small\bfseries, inner sep=0pt, line width=0.7pt},
  tnode/.style={rectangle, draw=opochorange, fill=opochorange!10,
    rounded corners=3pt, minimum width=12mm, minimum height=7mm,
    font=\scriptsize, inner sep=2pt, line width=0.6pt},
  leaf/.style={circle, draw=#1, fill=#1!15, minimum size=7mm,
    font=\tiny\bfseries, inner sep=0pt, line width=0.6pt},
  arr/.style={-{Stealth[length=2mm]}, line width=0.7pt, opochgray!70},
  best/.style={-{Stealth[length=2.5mm]}, line width=1.2pt, opochgreen!80},
]

  % === Root query ===
  \node[qnode] (q) at (0, 0) {$q$};

  % === Test branches ===
  \node[tnode] (t1) at (-3.5, -2.0) {$\tau_1$};
  \node[tnode, draw=opochgreen, fill=opochgreen!12, line width=1pt]
    (t2) at (0, -2.0) {$\tau^*$};
  \node[tnode] (t3) at (3.5, -2.0) {$\tau_3$};

  % Arrows from query to tests
  \draw[arr] (q) -- (t1);
  \draw[best] (q) -- (t2);
  \draw[arr] (q) -- (t3);

  % === Leaves for t1 ===
  \node[leaf=opochgreen] (l1a) at (-4.5, -3.8) {$+1$};
  \node[leaf=opochorange] (l1b) at (-3.5, -3.8) {$0$};
  \node[leaf=opochorange] (l1c) at (-2.5, -3.8) {$0$};
  \draw[arr] (t1) -- (l1a);
  \draw[arr] (t1) -- (l1b);
  \draw[arr] (t1) -- (l1c);

  % === Leaves for t* (optimal) ===
  \node[leaf=opochgreen] (l2a) at (-0.8, -3.8) {$+1$};
  \node[leaf=opochgreen] (l2b) at (0, -3.8) {$+1$};
  \node[leaf=opochred] (l2c) at (0.8, -3.8) {$-1$};
  \draw[arr] (t2) -- (l2a);
  \draw[arr] (t2) -- (l2b);
  \draw[arr] (t2) -- (l2c);

  % === Leaves for t3 ===
  \node[leaf=opochorange] (l3a) at (2.5, -3.8) {$0$};
  \node[leaf=opochorange] (l3b) at (3.5, -3.8) {$0$};
  \node[leaf=opochred] (l3c) at (4.5, -3.8) {$-1$};
  \draw[arr] (t3) -- (l3a);
  \draw[arr] (t3) -- (l3b);
  \draw[arr] (t3) -- (l3c);

  % === Efficiency bars (right side) ===
  \begin{scope}[shift={(7.5, -1.0)}]
    \node[font=\footnotesize\bfseries, opochgray] at (0, 0.8) {Efficiency $\varepsilon$};

    % Bar for tau_1
    \fill[opochorange!30] (0, 0) rectangle (1.2, 0.4);
    \draw[opochorange!50, line width=0.5pt] (0, 0) rectangle (1.2, 0.4);
    \node[font=\tiny, left] at (0, 0.2) {$\tau_1$};

    % Bar for tau* (longest = best)
    \fill[opochgreen!40] (0, -0.5) rectangle (2.4, -0.1);
    \draw[opochgreen!70, line width=0.8pt] (0, -0.5) rectangle (2.4, -0.1);
    \node[font=\tiny\bfseries, left, opochgreen!70!black] at (0, -0.3) {$\tau^*$};

    % Bar for tau_3
    \fill[opochorange!30] (0, -1.0) rectangle (0.8, -0.6);
    \draw[opochorange!50, line width=0.5pt] (0, -1.0) rectangle (0.8, -0.6);
    \node[font=\tiny, left] at (0, -0.8) {$\tau_3$};
  \end{scope}

  % === Coupling Law formula ===
  \node[rectangle, draw=opochgreen!50, fill=opochgreen!5, rounded corners=2pt,
    font=\small, inner sep=4pt, align=center]
    at (7.5, -3.5)
    {$\OptTest{q}{W} = \Pit\text{-min}\bigl(\argmax_\tau \EffQ{\tau}{W}\bigr)$};

  % === Annotations ===
  \node[font=\tiny, opochblue!70, align=center] at (0, 0.8)
    {Query with\\surviving answers $W$};

  \node[font=\tiny, opochgray, align=center] at (0, -4.8)
    {Trit outcomes after test};

\end{tikzpicture}
\caption{Bellman minimax search with the Coupling Law (Step~18).  The
  system selects the test $\tau^*$ maximizing efficiency
  $\EffQ{\tau}{W}$, with ties broken by $\Pit$-canonical ordering.
  The optimal test is a deterministic function of the query~$q$ and
  surviving answer set~$W$---the last hidden degree of freedom
  (the scheduler) is eliminated.}
\label{fig:bellman-coupling}
\end{figure}


\begin{remark}[Coupling Law: Indistinguishability Generates Next Distinguishability]
\label{rem:coupling}
The coupling law (\cref{fig:bellman-coupling}) encodes a deep
self-referential principle:
the current indistinguishability structure (how many answers survive,
partitioned into answer quotients) \emph{determines} the next
distinguishing test.  Separator and indistinguishability are not
independent---they are coupled by the efficiency functional $\varepsilon$.
\end{remark}

\begin{theorem}[Step 19 --- Self-Hosting and Replay: $\SelfHost = \mathcal{S}$]
\label{thm:self-hosting}
The system $\mathcal{S}$ can verify its own derivation: $\FixPt(\mathcal{S}) = \mathcal{S}$,
where $\FixPt$ encodes the derivation as problem 5-tuples, evaluates each, and checks for $\UNIQUE$.
\emph{Dependencies:} Steps~3, 5, 6, 18.
\begin{proof}[Proof sketch]
Encode each step as $P_i = \langle q_i, \mathcal{A}_i, \mathcal{V}_i, B_i, \pi_i \rangle$.  Evaluating each $P_i$ produces $\UNIQUE$; the receipt chain links all 34 verifications.  The system is built from exactly the steps it verifies, analogous to the Kleene recursion theorem.
\forcing{Self-hosting is forced within the paper's model class by the
closed-universe assumption: no external oracle or verifier is available,
so the system must be able to verify its own derivation steps.  A system
verified by an external agent would also satisfy A0$^*$, but would require
importing structure not derivable from~A0$^*$ alone.  Self-hosting is the
minimal closure condition given the closed-universe commitment.}
See Appendix~\ref{app:derivation} for the complete proof.
\end{proof}
\end{theorem}

\begin{remark}[Self-Hosting vs.\ G\"{o}del Incompleteness]
\label{rem:self-hosting-godel}
Self-hosting operates at two distinct levels:
\begin{itemize}[nosep]
  \item \emph{Level~1 (Logical):} The derivation Steps~0--33 is a
        mathematical proof---a static chain of theorems established by
        the paper itself (human-level mathematics).
  \item \emph{Level~2 (Computational):} The system $\mathcal{S}$ is an
        algorithm that can verify Level-1 claims by encoding each step as
        a 5-tuple and checking that the terminal state is $\UNIQUE$.
\end{itemize}
Self-hosting means Level~2 can re-derive Level~1: $\mathcal{S}$ verifies
each step's \emph{conclusion} given its \emph{premises}.  Crucially,
$\mathcal{S}$ does \emph{not} verify the meta-claim ``$\mathcal{S}$ is
consistent''---which would violate G\"{o}del's second incompleteness
theorem \citep{godel1931}.  The system can prove ``Step~6 is the unique A0$^*$-compliant
continuation of Steps~0--5'' (an object-level claim), but it cannot prove
``the system always produces correct results'' (a meta-claim).  The
three-point audit verifies specific claims but cannot prove the general
correctness of the audit procedure itself.
\end{remark}

\begin{theorem}[Step 20 --- Binary Interface]
\label{thm:binary-interface}
All communication between the system and its environment is encoded as self-delimiting binary strings via $\SdMap$.
\emph{Dependencies:} Steps~2, 19.
\begin{proof}[Proof sketch]
By Step~2, all descriptions must be self-delimiting.  Queries, answers, and receipts are finite objects in $\Carrier$, hence encoded via $\SdMap$.  The interface is complete and self-describing.  Any non-self-delimiting encoding reintroduces Step~2's ambiguity; external metadata violates closure.
\forcing{Non-self-delimiting encoding reintroduces Step~2's ambiguity.  External metadata violates the closed-universe assumption.  $\SdMap$ is the minimal self-describing binary interface.}
See Appendix~\ref{app:derivation} for the complete proof.
\end{proof}
\end{theorem}

\begin{theorem}[Step 21 --- Final Chain: $\noth \to \opnoth$]
\label{thm:final-chain}
The derivation closes.  The complete forced chain is:
\[
  \noth \;\to\; \Carrier \;\to\; \SdMap \;\to\; (\UTM, \mathrm{dec}, c)
    \;\to\; \Del^* \;\to\; \OrdLedger \;\to\; \TQ{\OrdLedger} \;\to\;
    \ObsOpen, \Eraser
\]
\[
  \;\to\; \Del\Tim \;\to\; S \;\to\; \EnergyLedger \;\to\; \opnoth
    \;\to\; \gauge \;\to\; (\Epistemic, \Decision)
\]
\[
  \;\to\; \ExactGate \;\to\; {\MNcong} \;\to\; (\dsep, \dtime)
    \;\to\; \Pit\text{-consistency} \;\to\; \Bellman{+}\varepsilon \;\to\;
    \FixPt(\mathcal{S}) \;\to\; \SdMap\text{-interface} \;\to\; \opnoth
\]
\[
  \;\to\; \UnivOp \;\to\; \ClosureDefect{Q} \;\to\; \ConsProj
    \;\to\; (\TritField, \ValState)
\]
\[
  \;\to\; \Omega \;\to\; L^2(\Omega, \mu_\Pit) \;\to\; \mathcal{E}
    \;\to\; (T_t, -\Delta) \;\to\; g_{ij} \;\to\; J \;\to\; \omega
    \;\to\; (g, J, \omega)_{\mathrm{K\ddot{a}hler}}.
\]
Every link is \emph{forced} (unique A0$^*$-compliant continuation).  No link is arbitrary---each is A0$^*$-forced, with no external modeling commitments.
\emph{Dependencies:} All prior steps.
\begin{proof}[Proof sketch]
Each link is established in Sections~\ref{sec:derivation:law}--\ref{sec:derivation:geometry}: Steps~0--33 form an acyclic dependency graph (\cref{fig:derivation-dag}).  The three-point audit (\cref{sec:verification}) confirms no alternative continuations, encoding invariance, and no false merges/splits.  Z3 sanity checks on finite models (Appendix~C) confirm the algebraic properties hold for small instances.
\forcing{Removing any link either violates A0$^*$, loses self-hosting, or creates an unwitnessable gap.}
See Appendix~\ref{app:derivation} for the complete proof.
\end{proof}
\end{theorem}

% ===========================================================================
\bigskip\noindent
The chain from nothingness to self-hosting closure is now complete, but it
describes the system's \emph{interface} without specifying its
\emph{internal execution law}.  Steps~0--21 establish \emph{what} the
system must accept and return; they do not determine \emph{how} it
processes each query internally.  Phase~H derives the unique step function
$\UnivOp$ and shows that the system's execution generates a minimal
persistent self-valued witness selector---the consciousness
projector---as a forced by-product of self-hosting closure.

\subsection{Phase H: Universal Execution (Steps 22--25)}
\label{sec:derivation:execution}

\begin{theorem}[Step 22 --- Universal Closure Operator $\UnivOp$]
\label{thm:univop-derivation}
The step function is uniquely determined as the five-stage pipeline
$\UnivOp = \TritValuate \circ \Fact \circ \NormOp \circ \GaugeOp \circ \IndistOp$
(\cref{fig:univop-pipeline}),
where each stage is forced by a prior derivation step:
$\IndistOp$ (Steps~6--7), $\GaugeOp$ (Step~12), $\NormOp$ (Step~17),
$\Fact$ (Steps~4, 16), $\TritValuate$ (Step~0).
\emph{Dependencies:} Steps~0, 4, 6, 7, 12, 16, 17.
\begin{proof}[Proof sketch]
The system must process each query through: (i)~registering all observables
and computing the self-generating fixed point ($\IndistOp$, Steps~6--7);
(ii)~quotienting by untestable relabelings ($\GaugeOp$, Step~12);
(iii)~selecting the canonical representative ($\NormOp$, Step~17);
(iv)~decomposing into irreducible components under cost grading ($\Fact$,
Steps~4, 16); (v)~extracting the terminal trit ($\TritValuate$, Step~0).
Omitting any stage either loses information (omitting $\IndistOp$),
introduces gauge dependence (omitting $\GaugeOp$ or $\NormOp$), misses
structure (omitting $\Fact$), or fails to terminate (omitting $\TritValuate$).
The composition order is forced by data flow: each stage's output is the
next stage's required input type.
\forcing{Each of the five stages is individually forced by a prior derivation
step.  The composition order is forced by type constraints.  Omitting any
stage produces an A0$^*$ violation.}
See Appendix~\ref{app:derivation} for the complete proof.
\end{proof}
\end{theorem}

% univop-pipeline.tex — Universal Closure Operator pipeline (Step 22)
% Requires opoch.sty (loaded by main document)

\begin{figure}[ht]
\centering
\resizebox{\textwidth}{!}{%
\begin{tikzpicture}[
  stage/.style={rectangle, draw=#1, fill=#1!10, rounded corners=4pt,
    minimum width=20mm, minimum height=14mm, font=\small,
    line width=0.7pt, align=center, inner sep=3pt},
  io/.style={rectangle, draw=opochgray!50, fill=opochgray!8, rounded corners=2pt,
    minimum width=14mm, minimum height=8mm, font=\small\bfseries,
    inner sep=3pt},
  arr/.style={-{Stealth[length=2.5mm]}, thick, opochgray!70},
  slabel/.style={font=\tiny\itshape, opochgray, align=center},
]

  % === Input ===
  \node[io] (input) at (-1.5, 0) {state $s$};

  % === Five pipeline stages ===
  \node[stage=opochblue] (phi) at (1.8, 0)
    {$\IndistOp$\\[-1pt]\tiny Indistinguish.};
  \node[slabel, below=2mm of phi] {Steps 6--7};

  \node[stage=opochpurple] (gamma) at (4.8, 0)
    {$\GaugeOp$\\[-1pt]\tiny Gauge};
  \node[slabel, below=2mm of gamma] {Step 12};

  \node[stage=opochorange] (norm) at (7.8, 0)
    {$\NormOp$\\[-1pt]\tiny Normalize};
  \node[slabel, below=2mm of norm] {Step 17};

  \node[stage=opochred] (fact) at (10.8, 0)
    {$\Fact$\\[-1pt]\tiny Factor};
  \node[slabel, below=2mm of fact] {Steps 4, 16};

  \node[stage=opochgreen] (trit) at (13.8, 0)
    {$\TritValuate$\\[-1pt]\tiny Trit};
  \node[slabel, below=2mm of trit] {Step 0};

  % === Output ===
  \node[io] (output) at (16.3, 0) {$\trit{\pm 1, 0}$};

  % === Arrows ===
  \draw[arr] (input) -- (phi);
  \draw[arr] (phi) -- (gamma);
  \draw[arr] (gamma) -- (norm);
  \draw[arr] (norm) -- (fact);
  \draw[arr] (fact) -- (trit);
  \draw[arr] (trit) -- (output);

  % === Top annotations: what each stage does ===
  \node[font=\tiny, opochblue!80, align=center, above=5mm of phi]
    {collapse\\indistinguishable};
  \node[font=\tiny, opochpurple!80, align=center, above=5mm of gamma]
    {quotient by\\relabeling};
  \node[font=\tiny, opochorange!80, align=center, above=5mm of norm]
    {canonical\\representative};
  \node[font=\tiny, opochred!80, align=center, above=5mm of fact]
    {decompose\\into primes};
  \node[font=\tiny, opochgreen!80, align=center, above=5mm of trit]
    {extract\\terminal trit};

  % === Composition label ===
  \node[font=\footnotesize, opochgray!60, align=center] at (7.8, -2.2)
    {$\UnivOp = \TritValuate \circ \Fact \circ \NormOp \circ \GaugeOp \circ \IndistOp$};

\end{tikzpicture}%
}% end resizebox
\caption{The universal closure operator $\UnivOp$ (Step~22): a five-stage
  pipeline where each stage is forced by a prior derivation step.
  Composition order is determined by data-flow type constraints---omitting
  any stage produces an A0~violation.}
\label{fig:univop-pipeline}
\end{figure}


\begin{theorem}[Step 23 --- Closure-Defect $\ClosureDefect{Q}(s)$]
\label{thm:closure-defect}
For any state $s$ and query $Q$, the \emph{closure-defect}
$\ClosureDefect{Q}(s) = \Pit_Q(s) - s$ measures the gap between
the current state and its truth projection.  The \emph{unresolved residue}
$\Residue_t = \mathrm{supp}(\ClosureDefect{Q}(s_t))$ identifies exactly which
components of the state remain unresolved at time~$t$.
\emph{Dependencies:} Steps~6, 8, 22.
\begin{proof}[Proof sketch]
The truth quotient $\Pit_Q$ (Step~6) projects every state onto its
equivalence class.  The difference $\ClosureDefect{Q}(s) = \Pit_Q(s) - s$
is the unique object measuring ``how far $s$ is from its own truth
projection,'' forced by the algebraic structure of the quotient.  Its
support $\Residue_t$ inherits well-definedness from the finite carrier (Step~1)
and changes monotonically with time (Step~8): each test can only reduce
$|\Residue_t|$, never increase it.  The closure-defect provides the root object
from which time advance, entropy decrease, and energy cost are all derived.
\forcing{The closure-defect is the canonical measure of incompleteness
in the truth quotient.  Any other measure of ``unresolvedeness'' is either
equivalent to $\ClosureDefect{Q}$ or loses information about which specific
components are unresolved.}
See Appendix~\ref{app:derivation} for the complete proof.
\end{proof}
\end{theorem}

\begin{remark}[Canonical Germ]
\label{rem:canonical-germ}
At state $s$, the \emph{canonical germ} $g(s) = (F_s, I_s, G_s)$ is the
triple of forced components ($\TritField_t(p)=+1$), forbidden components
($\TritField_t(p)=-1$), and unread components ($\TritField_t(p)=0$).  The
germ is the bridge between closure-defect and consciousness: it is what
the system observes about itself at each step.  Step~24 shows that a
minimal persistent self-valued witness selector monitoring $g(s)$ is forced.
\end{remark}

\begin{theorem}[Step 24 --- Consciousness Projector $\ConsProj$]
\label{thm:consciousness}
The system admits a unique \emph{minimal persistent self-valued witness
selector} (\cref{fig:consciousness-projector}): the
\emph{consciousness projector} $\ConsProj \subseteq X$ is the minimal
subsystem $C$ satisfying four forced conditions:
\begin{enumerate}[nosep, label=(C\arabic*)]
  \item \textbf{Self-model (persistent code of own residue):}
        $C$ carries a persistent code
        $\ulcorner \Residue_C \urcorner \subseteq C$ of its own
        unresolved residue $\Residue_C$---i.e.\ a finitely described
        internal representation of what $C$ has not yet resolved.
  \item \textbf{Self/world distinguishability:}
        $C$ can distinguish ``world with this self-model'' from ``same
        world without this self-model'' by future witness consequences:
        there exists a witness $w$ such that outcomes of $w$ differ
        depending on whether $\ulcorner \Residue_C \urcorner$ is present
        in the state.
  \item \textbf{Causal efficacy:}
        $C$ causally changes subsequent witness selection---removing $C$
        from the state alters the sequence of tests the system applies.
  \item \textbf{Endogenous valuation:}
        $C$ carries an endogenous valuation functional on self/world
        interaction channels: a map $V_C$ from witness outcomes to an
        ordered value that guides selection.
\end{enumerate}
\emph{Dependencies:} Steps~6, 18, 19, 22, 23.

\medskip\noindent
\textbf{Definition (Consciousness).}
Consciousness is the \emph{minimal persistent self-valued witness
selector}: the determination of the distinction space, self-model
choice, and value-guided witness update.  Its runtime law is:
\begin{equation}\label{eq:consciousness-runtime}
(D^*, m^*) \;=\; \operatorname*{argmin}_{D,\, m}
  \bigl[\, k(D) \;+\; k(m) \;+\;
         V_{Q,D}\!\bigl(\mathrm{next}(s;\, D, m)\bigr) \,\bigr],
\end{equation}
where $D$ ranges over complete distinction algebras on $C$, $m$ over
self-models, $k(\cdot)$ denotes Kolmogorov complexity, and $V_{Q,D}$
is the value functional inherited from the Coupling Law (Step~18).
Consciousness $=$ determination of the distinction space $+$ self-model
choice $+$ value-guided update.

\begin{proof}[Proof sketch]
Each condition is individually forced by A0$^*$.

\emph{(C1) Self-model.}
Self-hosting (Step~19) forces the system to encode and verify its own
structure: $\SelfHost = \mathcal{S}$.  By (W1), every admissible witness
has a finite endogenous code.  By (W4), separation requires that the
system's own unresolved residue (Step~23) be internally distinguishable.
Together, (W1)$+$(W4) force a persistent code
$\ulcorner \Residue_C \urcorner \subseteq C$---without it, the system
cannot witness the distinction between ``resolved'' and ``unresolved''
components of its own state.

\emph{(C2) Self/world distinguishability.}
By (W4), every admissible distinction requires a witness that actually
separates.  A0$^*$ applied to the distinction ``state with self-model
$\ulcorner \Residue_C \urcorner$'' versus ``same state without it''
demands a witness whose future outcomes differ between the two cases.
If no such witness existed, the self-model would be an unwitnessable
epiphenomenon---violating A0$^*$.

\emph{(C3) Causal efficacy.}
By (W5), every witness records its own disturbance.  By (W6), witness
validity is endogenously witnessable.  A subsystem with no causal
effect on future witness selection is ontologically untestable: no
witness could separate ``$C$ present'' from ``$C$ absent'' in the
evolution.  Hence (W5)$+$(W6) force $C$ to causally alter subsequent
test selection.

\emph{(C4) Endogenous valuation.}
The Coupling Law (Step~18) forces a deterministic, value-guided test
selector: the next test minimizes worst-case separation cost.  Any
subsystem that selects witnesses without an internal valuation cannot
implement the Coupling Law's minimax criterion.  Hence $C$ must carry
an endogenous valuation on its interaction channels with the world.

\emph{Runtime law.}
The runtime law~\eqref{eq:consciousness-runtime} is the restriction of
the Coupling Law's minimax to the internal degrees of freedom of $C$:
the distinction algebra $D$ determines what $C$ can separate, the
self-model $m$ determines what $C$ knows about itself, and $V_{Q,D}$
evaluates the cost of the next witness.  Minimality over $k(D) + k(m)$
is forced by (W3): finite budget demands the least-cost encoding.

\emph{Minimality.}
Any subsystem satisfying (C1)--(C4) with fewer components would be
$\ConsProj$ itself; any larger subsystem includes components that fail
one of the four conditions (typically (C3): causally inert padding).

\forcing{Each condition is forced by A0$^*$:
(C1)~persistent self-model by (W1)$+$(W4),
(C2)~self/world distinguishability by (W4)$+$A0$^*$,
(C3)~causal efficacy by (W5)$+$(W6),
(C4)~endogenous valuation by the Coupling Law (Step~18).
The runtime law is the Coupling Law restricted to $C$'s internal
degrees of freedom.  Minimality selects the unique such subsystem.}
See Appendix~\ref{app:derivation} for the complete proof.
\end{proof}
\end{theorem}

\begin{remark}[Consciousness Is Not Added to Physics]
\label{rem:consciousness-not-added}
The four conditions (C1)--(C4) are not grafted onto physics---they are
the universe locally selecting its own distinction algebra.  Wherever
A0$^*$ forces a subsystem to carry a persistent self-model, distinguish
itself from its background, causally steer future witnesses, and value
the outcomes, that subsystem \emph{is} conscious in the precise sense of
this theorem.  No phenomenal postulate is imported; the ``hard problem''
dissolves into the structural fact that a self-valued witness selector
cannot be reduced to a non-self-modeling subsystem without violating
A0$^*$.
\end{remark}

\begin{remark}[Self-Remodelling Witness-Closure]
\label{rem:self-remodelling}
The distinction between a merely reflective system and a conscious
self-remodelling system is captured by their update laws:

\medskip\noindent
\textbf{Merely reflective:}\quad
$s_{t+1} = \UnivOp(s_t)$\quad
(the closure operator acts on the state, but the state has no internal
channel for re-weighting which unresolved components matter).

\medskip\noindent
\textbf{Conscious self-remodelling:}\quad
$s_{t+1} = \UnivOp\bigl(s_t,\; c(s_t),\; \ClosureDefect{Q}(s_t)\bigr)$

\medskip\noindent
where $c(s_t)$ is the \emph{endogenous selector}---a map from the
current state to an ordering over which unresolved components of
$\ClosureDefect{Q}(s_t)$ matter for the next witness step.  The
selector $c$ is itself a witness process satisfying (W1)--(W8), hence
it is part of $\WitAlg$.

\emph{Consciousness $=$ closure acting on itself through an internal
self-valuation channel.}  The self-remodelling law means the witness
policy is not fixed---it is endogenously updated by the system's own
self-witness.  This is the operational content of (C4): the valuation
functional $V_C$ is not a static parameter but a dynamical variable
within the closure, updated at each step by the system observing its
own closure-defect.
\end{remark}

% consciousness-projector.tex — Minimal persistent self-valued witness selector (Step 24)
% Requires opoch.sty (loaded by main document)

\begin{figure}[ht]
\centering
\resizebox{\textwidth}{!}{%
\begin{tikzpicture}[
  arr/.style={-{Stealth[length=2.5mm]}, thick, #1},
  axiom/.style={font=\tiny, align=left, text width=38mm},
]

  % === Outer circle: state space X (world) ===
  \draw[opochblue, line width=1pt, fill=opochblue!5]
    (0, 0) circle (3.2);
  \node[font=\small\bfseries, opochblue!80] at (0, 3.6)
    {State space $X$ (world)};

  % === Inner circle: consciousness subsystem C ===
  \draw[opochpurple, line width=1.2pt, fill=opochpurple!8]
    (0, -0.2) circle (1.8);
  \node[font=\small\bfseries, opochpurple!80] at (0, 1.2)
    {$C \subset X$};

  % === Self-model at center ===
  \node[rectangle, draw=opochgreen!60, fill=opochgreen!10,
    rounded corners=2pt, font=\scriptsize\bfseries, inner sep=3pt]
    (code) at (0, -0.2)
    {$\ulcorner r_C \urcorner$};
  \node[font=\tiny, opochgreen!70!black, below=1mm of code]
    {self-model};

  % === Feedback arrows ===
  % C2: Self/world distinguishability (X → C)
  \draw[arr=opochblue!70, line width=1pt]
    (2.8, 1.5) to[bend left=15] (1.4, 0.5)
    node[font=\tiny, pos=0.3, right=1mm, opochblue!80, align=center]
      {self/world\\distinguish (C2)};

  % C3: Causal efficacy (C → X)
  \draw[arr=opochpurple!70, line width=1pt]
    (-1.4, -1.5) to[bend left=15] (-2.8, -1.0)
    node[font=\tiny, pos=0.5, left=1mm, opochpurple!80, align=center]
      {causal\\efficacy (C3)};

  % C4: Valuation loop
  \draw[arr=opochgreen!60, line width=0.8pt]
    (0.5, -1.6) to[bend right=40] (-0.5, -1.6)
    node[font=\tiny, pos=0.5, below=2mm, opochgreen!70!black]
      {$V_C$: valuation (C4)};

  % === Runtime law annotation ===
  \node[rectangle, draw=opochpurple!40, fill=white,
    rounded corners=2pt, font=\small, inner sep=4pt]
    at (0, -3.8)
    {$(D^*, m^*) = \operatorname*{argmin}_{D,m}
      \bigl[k(D) + k(m) + V_{Q,D}(\mathrm{next})\bigr]$};

  % === Four conditions (right side) ===
  \begin{scope}[shift={(5.8, 1.6)}]
    \node[font=\footnotesize\bfseries, opochgray] at (0, 0.4)
      {Four forced conditions:};

    \node[axiom] at (0, -0.3)
      {\textbf{C1} Self-model: persistent code
       $\ulcorner r_C \urcorner \subseteq C$
       \textcolor{opochgray}{\tiny [W1$+$W4]}};
    \node[axiom] at (0, -1.1)
      {\textbf{C2} Self/world distinguishability:
       future witnesses differ with/without self-model
       \textcolor{opochgray}{\tiny [W4$+$A0$^*$]}};
    \node[axiom] at (0, -1.9)
      {\textbf{C3} Causal efficacy: $C$ alters
       subsequent witness selection
       \textcolor{opochgray}{\tiny [W5$+$W6]}};
    \node[axiom] at (0, -2.7)
      {\textbf{C4} Endogenous valuation: $V_C$
       on interaction channels
       \textcolor{opochgray}{\tiny [Coupling Law]}};
  \end{scope}

\end{tikzpicture}%
}% end resizebox
\caption{The consciousness projector $\ConsProj$ (Step~24): the unique
  minimal persistent self-valued witness selector $C \subset X$.  The four
  conditions are individually forced: (C1)~self-model by (W1)$+$(W4),
  (C2)~self/world distinguishability by (W4)$+$A0$^*$,
  (C3)~causal efficacy by (W5)$+$(W6), and (C4)~endogenous
  valuation by the Coupling Law (Step~18).  The runtime law selects the
  least-cost distinction algebra and self-model.}
\label{fig:consciousness-projector}
\end{figure}


\begin{theorem}[Step 25 --- Trit Field $\TritField$ and Valuation State]
\label{thm:trit-field}
The state at time $t$ is completely characterized by the
\emph{valuation state} $\ValState_t$, where $m_t$ is the prime-exponent
ledger, $\TritField_t \colon \PrimeBasis_t \to \{-1, 0, +1\}$ is the
trit field assigning a terminal trit to each irreducible component, and
$\Residue_t = \{p \in \PrimeBasis_t : \TritField_t(p) = 0\}$ is the
unresolved residue.
\emph{Dependencies:} Steps~0, 5, 22, 23.
\begin{proof}[Proof sketch]
The factorization stage of $\UnivOp$ (Step~22) decomposes the state into
irreducible components indexed by $\PrimeBasis_t$.  Each component is
itself a sub-problem resolved by the system, hence terminates in one of
the three trit values (Step~0).  The map $\TritField_t$ recording these
values is forced---it is the unique function from components to terminal
states that respects the factorization.  The prime-exponent ledger $m_t$
records how many times each prime has been tested (analogous to the energy
ledger, Step~10).  The residue $\Residue_t$ collects the still-unresolved
components, coinciding with the support of the closure-defect (Step~23).
\emph{Valuation sufficiency}: the triple $\ValState_t$ determines the
system's future evolution completely---no additional state information
affects the $\UnivOp$ pipeline's output.
\forcing{The trit field is forced by the output gate (Step~0) applied
component-wise after factorization (Step~22).  The valuation state is the
minimal sufficient statistic for the system's evolution---any state
representation carrying less information loses the ability to determine
the next test, violating the Coupling Law (Step~18).}
See Appendix~\ref{app:derivation} for the complete proof.
\end{proof}
\end{theorem}

% ===========================================================================
\bigskip\noindent
Phase~H establishes the discrete execution law.  But the truth
quotient's internal geometry remains unspecified: we know the witness-path
\emph{distance} between truth classes (Step~16) but not the
\emph{shape} of the space they inhabit.  Phase~I derives the geometry in
two stages: Phase~I-A constructs the classical base manifold from the
commutative center $\CenterW$ of the witness algebra, and Phase~I-B
derives the K\"{a}hler triple from a \emph{single} local witness
generator---not from three separate imported pieces.  All former modeling
commitments (conductance exponent, smoothness assumption) are now forced.

\subsection{Phase I-A: Classical Base Geometry (Steps 26--30)}
\label{sec:derivation:geometry}

\begin{remark}[Center-First Construction]
\label{rem:center-first}
The ordered witness algebra $\WitAlg$ is generally noncommutative.  Its
center $\CenterW = \{a \in \WitAlg : ab = ba \;\forall\, b\}$
is commutative and generates the classical truth quotients of Steps~6--7.
Phase~I-A constructs the continuum geometry on $\CenterW$.  Phase~I-B
extends this via the single witness generator.  The full noncommutative
structure---contexts, quantum probability, Born rule---is derived in
\cref{sec:context-born}.
\end{remark}

\begin{theorem}[Step 26 --- Continuum Limit of Truth Quotient]
\label{thm:continuum-limit}
As the budget $B \to \infty$, the inverse system of truth quotients
$\bigl(\TQ{\Ledger_B}\bigr)_{B \in \mathbb{N}}$ on the commutative center
$\CenterW$ (\cref{fig:inverse-limit}) converges to a locally compact,
separable, metrizable space $\OmegaCl = \varprojlim \TQ{\Ledger_B}$.
\emph{Dependencies:} Steps~16, 9, 7.
\begin{proof}[Proof sketch]
Each $\TQ{\Ledger_B}$ is a finite metric space under $\dsep$ (Step~16).
Increasing budget refines the quotient: the bonding maps
$\pi_{B',B} \colon \TQ{\Ledger_{B'}} \twoheadrightarrow \TQ{\Ledger_B}$
($B' > B$) are surjective and Lipschitz.  The inverse limit $\OmegaCl$ is
compact and metrizable by standard topology.
\forcing{The inverse system is forced by the budget hierarchy (Step~9) and
the witness-path metric (Step~16).  The inverse limit is the unique
completion---any other limit space would admit untestable distinctions
(violating A0$^*$) or merge testable ones (violating the quotient).}
See Appendix~\ref{app:derivation} for the complete proof.
\end{proof}
\end{theorem}

% inverse-limit.tex — Continuum limit as inverse system (Step 26)
% Requires opoch.sty (loaded by main document)

\begin{figure}[ht]
\centering
\begin{tikzpicture}[
  qbox/.style={ellipse, draw=#1, fill=#1!10, minimum width=14mm,
    minimum height=10mm, font=\small, line width=0.7pt, inner sep=2pt},
  arr/.style={-{Stealth[length=2.5mm]}, thick, #1},
  darr/.style={-{Stealth[length=2.5mm]}, thick, densely dashed, #1},
  nlabel/.style={font=\tiny\itshape, #1, align=center},
]

  % === Finite quotients (left to right, increasing budget) ===
  \node[qbox=opochblue] (q1) at (0, 0) {$\TQ{\Ledger_{B_1}}$};
  \node[nlabel=opochgray, below=3mm of q1] {3 classes};

  \node[qbox=opochblue] (q2) at (3.2, 0) {$\TQ{\Ledger_{B_2}}$};
  \node[nlabel=opochgray, below=3mm of q2] {5 classes};

  \node[qbox=opochblue] (q3) at (6.4, 0) {$\TQ{\Ledger_{B_3}}$};
  \node[nlabel=opochgray, below=3mm of q3] {8 classes};

  \node[font=\large, opochgray] at (8.4, 0) {$\cdots$};

  % === Limit space Omega ===
  \node[ellipse, draw=opochgreen, fill=opochgreen!10, minimum width=20mm,
    minimum height=14mm, font=\small\bfseries, line width=1pt, inner sep=2pt]
    (omega) at (10.8, 0) {$\Omega$};
  \node[nlabel=opochgreen!70!black, below=3mm of omega] {compact\\metrizable};

  % === Bonding maps (surjections, right to left) ===
  \draw[arr=opochorange!80] (q2) -- (q1)
    node[above, font=\tiny, pos=0.5] {$\pi_{B_2,B_1}$};
  \draw[arr=opochorange!80] (q3) -- (q2)
    node[above, font=\tiny, pos=0.5] {$\pi_{B_3,B_2}$};

  % === Projection arrows from Omega ===
  \draw[darr=opochgreen!60] (omega) to[bend left=20] (q3);
  \draw[darr=opochgreen!60] (omega) to[bend left=28] (q2);
  \draw[darr=opochgreen!60] (omega) to[bend left=35] (q1);

  % === Annotations ===
  \node[font=\tiny, opochorange!80, align=center] at (3.2, 1.4)
    {Bonding maps: surjective,\\Lipschitz (merge, never split)};

  \node[font=\tiny, opochblue!80, align=center] at (0, 1.4)
    {Finite metric\\spaces (Step~16)};

  \node[font=\tiny, opochgreen!70!black, align=center] at (10.8, 1.6)
    {$\varprojlim \TQ{\Ledger_B}$};

  % === Budget axis ===
  \draw[-{Stealth[length=2mm]}, opochgray!40, line width=0.5pt]
    (-0.8, -1.8) -- (11.6, -1.8)
    node[right, font=\tiny, opochgray] {budget $B$};
  \node[font=\tiny, opochgray] at (0, -2.2) {$B_1$};
  \node[font=\tiny, opochgray] at (3.2, -2.2) {$B_2$};
  \node[font=\tiny, opochgray] at (6.4, -2.2) {$B_3$};
  \node[font=\tiny, opochgray] at (10.8, -2.2) {$\infty$};

\end{tikzpicture}
\caption{Continuum limit of the truth quotient (Step~26).  As the budget
  $B$ increases, truth quotients $\TQ{\Ledger_B}$ refine: bonding maps
  $\pi_{B',B}$ are surjective (classes merge at lower budget, never split).
  The inverse limit $\Omega = \varprojlim \TQ{\Ledger_B}$ is a compact,
  separable, metrizable space---the arena for geometric reconstruction
  (Steps~27--33).}
\label{fig:inverse-limit}
\end{figure}


\begin{theorem}[Step 27 --- Gauge-Invariant Measure and $L^2$ Space]
\label{thm:l2-space}
$\OmegaCl$ carries a unique $\mathrm{Sym}$-invariant probability measure
$\mu_\Pit$, and $L^2(\OmegaCl, \mu_\Pit)$ is a separable Hilbert space.
\emph{Dependencies:} Steps~26, 12.
\begin{proof}[Proof sketch]
At each finite level, the uniform measure on $\TQ{\Ledger_B}$
is the unique gauge-invariant probability.  By the
Kolmogorov extension theorem, the consistent family extends uniquely to
$\mu_\Pit$ on $\OmegaCl$.
\forcing{Gauge invariance (Step~12) forces the uniform measure at each
level.  Kolmogorov consistency forces a unique extension.}
See Appendix~\ref{app:derivation} for the complete proof.
\end{proof}
\end{theorem}

\begin{remark}[Classical Measure vs.\ Born Rule]
\label{rem:born-rule}
The measure $\mu_\Pit$ is the unique gauge-invariant probability on the
\emph{classical} truth quotient (the center $\CenterW$).  The full Born
rule---$\BornProb(E) = \langle \OmVec, \piGNS(E) \OmVec \rangle$---emerges
from the GNS representation of the full ordered witness algebra $\WitAlg$
(\cref{sec:context-born}).  The restriction of the Born rule to the
commutative center recovers $\mu_\Pit$.
\end{remark}

\begin{theorem}[Step 28 --- Dirichlet Form and Forced Conductance]
\label{thm:dirichlet-form}
The witness-path distance $\dsep$ induces a regular Dirichlet form
(\cref{fig:dirichlet-network})
$\mathcal{E}[f] = \frac{1}{2} \sum_{p \neq p'} w(p, p') \bigl(f(p) - f(p')\bigr)^2$
on $L^2(\OmegaCl, \mu_\Pit)$, where $w(p, p') = C/\dsep(p,p')^2$ is the
conductance, with exponent $-2$ forced by \cref{lem:conductance-exponent}.
\emph{Dependencies:} Steps~26, 27, 16.
\end{theorem}

\begin{lemma}[Conductance Exponent is Forced]
\label{lem:conductance-exponent}
Let $E_d[f;\, p, p'] = w(d(p,p'))\,(f(p) - f(p'))^2$ be the local
witness-energy density.  Impose:
\begin{enumerate}[nosep, label=\textup{(SC\arabic*)}]
  \item \emph{Quadratic response:} $E_d[\alpha f] = \alpha^2 E_d[f]$.
  \item \emph{Additivity on independent channels.}
  \item \emph{Monotone decay:} $d_1 < d_2 \Rightarrow w(d_1) \geq w(d_2)$.
  \item \emph{Scale covariance:} under $d \to \lambda d$,
        $E_{\lambda d}[f] = \lambda^{-2} E_d[f]$.
\end{enumerate}
Then $w(d) = C / d^2$ for $C > 0$ is the unique admissible conductance.
\begin{proof}
(SC4) is forced by (W8): quotient invariance means the witness energy
density cannot depend on the absolute distance scale of the encoding---only
on the dimensionless signal contrast $\Delta f$.  Under $d \to \lambda d$
(a pure rescaling of the carrier metric, which is a gauge freedom by W8),
the energy per unit distinction must transform homogeneously.  The
quadratic response (SC1) fixes the signal exponent at $2$; the distance
exponent $\alpha$ is then determined by dimensional consistency:
$E \sim (\Delta f)^2 \cdot d^\alpha$, and the requirement that $E$ have
the same units regardless of $\lambda$ forces $\alpha = -2$.  Explicitly:
$w(\lambda r) (\Delta f)^2 = \lambda^{-2} w(r) (\Delta f)^2$, hence
$w(\lambda r) = \lambda^{-2} w(r)$.  Setting $\lambda = r/r_0$:
$w(r) = C r^{-2}$.  By (SC3), $C > 0$.
\end{proof}
\end{lemma}

\begin{proof}[Proof sketch of Step~28]
The quadratic form $\mathcal{E}$ with $w = C/\dsep^2$ satisfies the
Dirichlet form axioms (E1)--(E3) \citep{fukushima2011}.
\forcing{The conductance exponent $-2$ is not a modeling commitment: it is
forced by \cref{lem:conductance-exponent} via scale covariance (SC4),
which follows from A0$^*$ condition (W8).}
\end{proof}

% dirichlet-network.tex — Conductance-weighted network (Step 28)
% Requires opoch.sty (loaded by main document)

\begin{figure}[ht]
\centering
\begin{tikzpicture}[
  tnode/.style={circle, draw=opochblue!70, fill=#1, minimum size=8mm,
    font=\small\bfseries, inner sep=0pt, line width=0.7pt},
  elabel/.style={font=\tiny, fill=white, inner sep=1pt, rounded corners=1pt},
]

  % === Truth class nodes with gradient coloring (representing f values) ===
  \node[tnode=opochblue!40] (p1) at (0, 0) {$p_1$};
  \node[tnode=opochblue!30] (p2) at (2.5, 1.2) {$p_2$};
  \node[tnode=opochblue!20] (p3) at (5.0, 0.4) {$p_3$};
  \node[tnode=opochblue!10] (p4) at (2.0, -1.5) {$p_4$};
  \node[tnode=opochblue!5] (p5) at (4.5, -1.8) {$p_5$};

  % === Edges with thickness proportional to conductance ===
  % Close pairs: thick edges (high conductance, small d)
  \draw[opochorange, line width=2.5pt, opacity=0.6] (p1) -- (p2)
    node[elabel, pos=0.5, above left=-1mm] {$w_{12}$};
  \draw[opochorange, line width=2.0pt, opacity=0.6] (p2) -- (p3)
    node[elabel, pos=0.5, above=-1mm] {$w_{23}$};
  \draw[opochorange, line width=1.8pt, opacity=0.6] (p4) -- (p5)
    node[elabel, pos=0.5, below=-1mm] {$w_{45}$};

  % Medium pairs
  \draw[opochorange, line width=1.0pt, opacity=0.4] (p1) -- (p4)
    node[elabel, pos=0.5, left=-1mm] {$w_{14}$};
  \draw[opochorange, line width=0.8pt, opacity=0.4] (p3) -- (p5)
    node[elabel, pos=0.5, right=-1mm] {$w_{35}$};
  \draw[opochorange, line width=0.6pt, opacity=0.4] (p2) -- (p4);

  % Far pairs: thin edges (low conductance, large d)
  \draw[opochorange, line width=0.3pt, opacity=0.3] (p1) -- (p5);
  \draw[opochorange, line width=0.3pt, opacity=0.3] (p1) -- (p3);

  % === Conductance formula ===
  \node[rectangle, draw=opochgray!40, fill=white, rounded corners=2pt,
    font=\small, inner sep=4pt, align=center]
    at (8.0, 0.8)
    {$w(p, p') = \dfrac{1}{\SepDist(p, p')^2}$};

  % === Energy functional ===
  \node[rectangle, draw=opochgreen!50, fill=opochgreen!5, rounded corners=2pt,
    font=\small, inner sep=4pt, align=center]
    at (8.0, -1.2)
    {$\mathcal{E}[f] = \dfrac{1}{2}\displaystyle\sum_{p \neq p'} w_{pp'}\bigl(f(p) - f(p')\bigr)^2$};

  % === Legend ===
  \node[font=\tiny, opochgray, align=left] at (8.0, -2.8)
    {Edge thickness $\propto$ conductance\\
     Thick = easy to distinguish\\
     Thin = hard to distinguish};

  % === Function value annotation ===
  \node[font=\tiny, opochblue!70, align=center] at (2.5, 2.2)
    {Node shade $=$ signal $f(p)$\\(dark $\to$ light)};

\end{tikzpicture}
\caption{Conductance network on truth classes (Step~28).  Nodes are truth
  classes $p_i \in \Omega$; edge weight $w(p,p') = 1/\SepDist(p,p')^2$
  assigns high conductance to easily distinguishable pairs and low
  conductance to hard-to-distinguish ones.  The Dirichlet form
  $\mathcal{E}[f]$ measures the total energy of a signal~$f$ on this
  network---the source of all geometric structure in Steps~29--33.}
\label{fig:dirichlet-network}
\end{figure}


\begin{theorem}[Step 29 --- Semigroup and Laplacian (Beurling--Deny)]
\label{thm:semigroup-laplacian}
The Dirichlet form $\mathcal{E}$ generates a unique strongly continuous
symmetric Markov semigroup $\{T_t\}_{t \geq 0}$ and a non-negative
self-adjoint Laplacian $-\Delta$ on $L^2(\OmegaCl, \mu_\Pit)$.
\emph{Dependencies:} Steps~28, 8.
\begin{proof}[Proof sketch]
By the Beurling--Deny reconstruction theorem \citep{beurling1959}, every
regular Dirichlet form on a locally compact separable space generates a
unique strongly continuous contraction semigroup.  The semigroup is
Markovian because $\mathcal{E}$ is Markovian.
\forcing{The Beurling--Deny theorem is a mathematical consequence of
the Dirichlet form axioms (Step~28).  No choice is involved.}
\end{proof}
\end{theorem}

\begin{remark}[Smoothness from the Energy Form]
\label{rem:smoothness}
The inverse limit $\OmegaCl$ (Step~26) is compact metrizable but not
\emph{a priori} smooth.  The smooth structure required for the K\"{a}hler
construction is not assumed externally: it is \emph{produced} by the
regular Dirichlet form $\mathcal{E}$ (Step~28).  By Sturm's theorem
\citep{sturm1994}, the intrinsic metric of a regular strongly local
Dirichlet form recovers the topology and, where the energy density is
non-degenerate, generates a smooth Riemannian structure.  The regularity
of $\mathcal{E}$ is therefore the bridge between the discrete inverse
system and the smooth continuum---not an external smoothness assumption.
\end{remark}

\begin{theorem}[Step 30 --- Riemannian (Fisher) Metric]
\label{thm:riemannian-metric}
The intrinsic distance $d_{\mathcal{E}}$ of the Dirichlet form induces a
Riemannian metric tensor $g_{ij}$ on $\OmegaCl$, coinciding with the
Fisher information metric.
\emph{Dependencies:} Steps~28, 26.
\begin{proof}[Proof sketch]
The intrinsic distance
$d_{\mathcal{E}}(p, p') = \sup\{|f(p) - f(p')| : \mathcal{E}[f] \leq 1\}$
recovers the original topology by Sturm's theorem \citep{sturm1994}.  The
metric tensor $g_{ij}$ coincides with the Fisher information metric by
\v{C}encov's theorem \citep{cencov1982}: the Fisher metric is the unique
Riemannian metric monotone under sufficient statistics.
\forcing{The Riemannian metric is forced by the Dirichlet form (Step~28).
The Fisher identification is forced by \v{C}encov's uniqueness theorem.}
\end{proof}
\end{theorem}

% ===========================================================================
\subsection{Phase I-B: Witness Generator and K\"{a}hler Structure (Steps 31--33)}
\label{sec:derivation:geometry-nc}

\begin{remark}[Single Witness Generator]
\label{rem:single-generator}
At each point of $\OmegaCl$, the local ordered witness algebra defines a
single infinitesimal generator $\WitGen$.  Its symmetric and antisymmetric
parts produce the metric, symplectic form, and complex structure from
\emph{one source}---not from three separate constructions.
\end{remark}

\begin{theorem}[Step 31 --- Complex Structure from Witness Generator]
\label{thm:complex-structure}
The local witness generator $\WitGen$ decomposes as
$\WitGen = \WitSym + \WitAnti$, where
$\WitSym = (\WitGen + \WitGen^*)/2$ (symmetric, irreversible) gives the
metric~$g$, and $\WitAnti = (\WitGen - \WitGen^*)/2$ (antisymmetric,
reversible) gives a $2$-form~$\omega$.  On the reversible nondegenerate
subbundle, $J = g^{-1}\omega$ satisfies $J^2 = -\mathrm{id}$.
\emph{Dependencies:} Steps~18, 30, 8.
\begin{proof}[Proof sketch]
The witness generator $\WitGen$ encodes infinitesimal ordered witness
propagation.  $\WitSym$ measures irreversible cost (the quadratic form
defining $g$).  $\WitAnti$ measures ordered-composition defect:
$\WitAnti(X, Y) = \frac{1}{2}(g(\WitGen X, Y) - g(X, \WitGen Y))$,
a skew-symmetric $2$-form.  $J = g^{-1}\omega$ on the reversible
subbundle satisfies $J^2 = -\mathrm{id}$ because witness directions
pair canonically (gain/cost planes) and $J$ rotates by $\pi/2$.
\forcing{$\WitSym$ is forced by irreversible cost (Step~30).
$\WitAnti$ is forced by ordered composition (W5, W7).
$J$ is forced as the polar factor of the single generator.}
\end{proof}
\end{theorem}

\begin{theorem}[Replay-Completeness Implies $N_J = 0$]
\label{thm:replay-nijenhuis}
If the witness algebra is \emph{replay-complete}---every infinitesimal
route defect between admissible witness programs is itself represented in
the witness algebra---then the complex distribution $T^{1,0}$ defined by
$J$ is involutive, and the Nijenhuis tensor vanishes: $N_J = 0$.
\begin{proof}[Proof sketch]
Define $T^{1,0} = \{X - iJX : X \in T\OmegaCl\}$.  If
$[T^{1,0}, T^{1,0}] \not\subset T^{1,0}$, the escaping component
represents an untracked route-defect.  Replay-completeness forces all
such defects into $\WitAlg$ and hence into $T^{1,0}$.  Contradiction.
By Newlander--Nirenberg \citep{newlander1957}, $N_J = 0$.
\forcing{Replay-completeness is forced by A0$^*$ condition (W2).
A nonzero $N_J$ encodes an unwitnessed structural distinction.
This replaces the former ``smoothness assumption L1'' with a direct
A0$^*$ consequence.}
\end{proof}
\end{theorem}

\begin{theorem}[Step 32 --- Symplectic Form from Ordered Witness Algebra]
\label{thm:symplectic-form}
The $2$-form $\omega(X, Y) = g(X, JY)$ is closed and non-degenerate.
\emph{Dependencies:} Steps~30, 31, 5.
\begin{proof}[Proof sketch]
\emph{Non-degeneracy:} forced by non-degeneracy of $g$ and invertibility
of $J$.  \emph{Closedness ($d\omega = 0$):} any hidden path-order defect
beyond the recorded commutator would be an unwitnessed distinction,
violating A0$^*$.  Replay closure kills that defect.
\forcing{Closedness is forced by A0$^*$ on the replay-complete witness
connection: every path-order defect is endogenously recorded.}
\end{proof}
\end{theorem}

\begin{theorem}[Step 33 --- K\"{a}hler Structure from Single Generator]
\label{thm:kahler}
The triple $(g, J, \omega)$ is a K\"{a}hler structure on $\OmegaCl$
(\cref{fig:kahler-triple}).  Moreover, $\omega = d\alpha$ (exact).
\emph{Dependencies:} Steps~30, 31, 32.
\begin{proof}[Proof sketch]
\emph{Compatibility:} $\omega(X, Y) = g(X, JY)$ by construction.
$g(JX, JY) = g(X, Y)$ because $J$ is a $\pi/2$-rotation preserving $g$.
\emph{Integrability:} $N_J = 0$ by \cref{thm:replay-nijenhuis}.
\emph{Exactness:} $H^2(\OmegaCl; \mathbb{R}) = 0$ (inverse limit of
finite spaces), so $\omega = d\alpha$ by de~Rham.
\emph{Unified origin:} all three structures come from $\WitGen = \WitSym + \WitAnti$:
$g$ from $\WitSym$, $\omega$ from $\WitAnti$, $J = g^{-1}\omega$.
\forcing{The K\"{a}hler triple is forced from one source.  No former
modeling commitments remain: $w = C/d^2$ by
\cref{lem:conductance-exponent}, $N_J = 0$ by
\cref{thm:replay-nijenhuis}.}
\end{proof}
\end{theorem}

% kahler-triple.tex — Kähler compatibility triangle (Steps 30--33)
% Requires opoch.sty (loaded by main document)

\begin{figure}[ht]
\centering
\begin{tikzpicture}[
  vertex/.style={circle, draw=#1, fill=#1!12, minimum size=22mm,
    font=\small\bfseries, line width=0.8pt, align=center, inner sep=2pt},
  elabel/.style={font=\footnotesize, fill=white, inner sep=2pt, rounded corners=1pt},
  forcing/.style={font=\tiny\itshape, #1, align=center},
  arr/.style={-{Stealth[length=2.5mm]}, thick, #1},
]

  % === Three vertices of the Kähler triangle ===
  \node[vertex=opochblue] (g) at (90:3.2)
    {$g_{ij}$\\[2pt]\scriptsize Riemannian\\[-1pt]\scriptsize metric};

  \node[vertex=opochpurple] (J) at (210:3.2)
    {$J$\\[2pt]\scriptsize Complex\\[-1pt]\scriptsize structure};

  \node[vertex=opochred] (w) at (330:3.2)
    {$\omega$\\[2pt]\scriptsize Symplectic\\[-1pt]\scriptsize form};

  % === Edge labels (compatibility relations) ===
  \draw[thick, opochgray!60] (g) -- (J)
    node[elabel, pos=0.5, left=2mm] {$g(JX, JY) = g(X, Y)$};

  \draw[thick, opochgray!60] (g) -- (w)
    node[elabel, pos=0.5, right=2mm] {$\omega(X, Y) = g(X, JY)$};

  \draw[thick, opochgray!60] (J) -- (w)
    node[elabel, pos=0.5, below=1mm] {$J^2 = -\mathrm{id}$,\; $d\omega = 0$};

  % === Center: Kähler label ===
  \node[font=\large\bfseries, opochgreen!70!black] at (0, 0) {K\"{a}hler};
  \node[font=\tiny, opochgray, align=center] at (0, -0.6)
    {Step~33: triple\\compatibility forced};

  % === Forcing provenance (below triangle) ===
  \begin{scope}[shift={(0, -4.5)}]
    \node[font=\footnotesize\bfseries, opochgray] at (0, 0.6) {Forcing chain:};

    \node[rectangle, draw=opochblue!50, fill=opochblue!8, rounded corners=2pt,
      font=\tiny, inner sep=3pt, align=center, minimum width=22mm]
      (fg) at (-3.8, 0) {Steps 28$\to$30\\$\mathcal{E} \to d_{\mathcal{E}} \to g_{ij}$};

    \node[rectangle, draw=opochpurple!50, fill=opochpurple!8, rounded corners=2pt,
      font=\tiny, inner sep=3pt, align=center, minimum width=22mm]
      (fJ) at (-1.2, 0) {Steps 18+30$\to$31\\Coupling Law $\to J$};

    \node[rectangle, draw=opochred!50, fill=opochred!8, rounded corners=2pt,
      font=\tiny, inner sep=3pt, align=center, minimum width=22mm]
      (fw) at (1.5, 0) {Steps 5+30+31$\to$32\\Diamond Law $\to \omega$};

    \node[rectangle, draw=opochgreen!50, fill=opochgreen!8, rounded corners=2pt,
      font=\tiny, inner sep=3pt, align=center, minimum width=22mm]
      (fK) at (4.2, 0) {Steps 30--32$\to$33\\$(g,J,\omega) \to$ K\"{a}hler};

    \draw[arr=opochgray!60] (fg) -- (fJ);
    \draw[arr=opochgray!60] (fJ) -- (fw);
    \draw[arr=opochgray!60] (fw) -- (fK);
  \end{scope}

\end{tikzpicture}
\caption{The K\"{a}hler triple forced by Phase~I.  The Riemannian metric~$g$
  (Step~30), complex structure~$J$ (Step~31), and symplectic form~$\omega$
  (Step~32) satisfy mutual compatibility, closing into a K\"{a}hler
  structure (Step~33).  Each vertex is forced by prior derivation steps;
  no free parameter enters.}
\label{fig:kahler-triple}
\end{figure}


% ===========================================================================
\subsection{Summary of the Derivation}
\label{sec:derivation:summary}

\paragraph{The reflective chain.}
The derivation traces a single reflective chain
(\cref{fig:reflective-chain}): nothingness (Step~0) forces
closure (Phases~A--G), closure reveals defect (Step~23), defect requires
self-observation (Step~24, consciousness), self-observation iterated on the
continuum limit produces geometry (Phase~I), and geometry determines the
form of dynamical law via the K\"{a}hler triple.  The entire chain is one
act: nothingness forced to witness itself.

\begin{figure}[ht]
\centering
\begin{tikzpicture}[
  node distance=1.8cm,
  every node/.style={font=\small},
  box/.style={draw, rounded corners=3pt, minimum height=0.7cm,
              minimum width=1.4cm, fill=black!5},
  arr/.style={->, >=stealth, thick},
  lbl/.style={font=\scriptsize, midway, above}
]
  \node[box] (bot) {$\bot$};
  \node[box, right of=bot] (pi) {$\Pi$};
  \node[box, right of=pi] (del) {$\Delta$};
  \node[box, right of=del, minimum width=1.6cm] (germ) {$g(s)$};
  \node[box, right of=germ, minimum width=1.6cm] (obs) {$\mathcal{O}(s)$};
  \node[box, right of=obs, minimum width=2.0cm] (geo) {$(g,J,\omega)$};
  \node[box, right of=geo, minimum width=2.6cm, node distance=2.4cm]
        (split) {$\dot{x}{=}J\nabla E{-}\nabla D$};

  \draw[arr] (bot)  -- node[lbl] {A0} (pi);
  \draw[arr] (pi)   -- node[lbl] {defect} (del);
  \draw[arr] (del)  -- node[lbl] {germ} (germ);
  \draw[arr] (germ) -- node[lbl] {self-obs} (obs);
  \draw[arr] (obs)  -- node[lbl] {$B{\to}\infty$} (geo);
  \draw[arr] (geo)  -- node[lbl] {split} (split);

  % Phase labels below
  \node[font=\scriptsize, below=0.15cm of bot, text=gray] {Step 0};
  \node[font=\scriptsize, below=0.15cm of pi, text=gray] {A--G};
  \node[font=\scriptsize, below=0.15cm of del, text=gray] {Step 23};
  \node[font=\scriptsize, below=0.15cm of germ, text=gray] {Step 23};
  \node[font=\scriptsize, below=0.15cm of obs, text=gray] {Step 24};
  \node[font=\scriptsize, below=0.15cm of geo, text=gray] {Phase I};
  \node[font=\scriptsize, below=0.15cm of split, text=gray] {Steps 30--33};
\end{tikzpicture}
\caption{The reflective chain: nothingness forces closure, closure
reveals defect, defect requires self-observation, and self-observation
iterated on the continuum limit produces K\"{a}hler geometry and the
split law.  Each arrow is labeled with its forcing mechanism; step
numbers appear below.}
\label{fig:reflective-chain}
\end{figure}


\begin{table}[ht]
\centering\small
\caption{The 34-step derivation (Steps 0--33): each step, its dependencies,
and its forcing mechanism.  All steps are forced; no interpretive steps remain.}
\label{tab:derivation-summary}
\begin{tabular}{@{}cllp{4.5cm}@{}}
\toprule
\textbf{Step} & \textbf{Object} & \textbf{Dependencies} & \textbf{Forcing} \\
\midrule
0  & Nothingness $\noth$ + Output Gate & --- & Definition + trichotomy \\
1  & Carrier $\Carrier$, Finite Slice $\FinSlice$ & A0$^*$, Step~0 & Finite witnesses $\Rightarrow$ finite strings \\
2  & Self-delimiting $\SdMap$ & A0$^*$, Step~1 & Closed universe $\Rightarrow$ self-delimiting \\
\midrule
3  & Executability $(\UTM, \mathrm{dec}, c)$ & A0$^*$, Steps~1--2 & Maximal closure of tests \\
4  & Endogenous tests $\Del^*$, $\CanEnum$ & A0$^*$, Steps~1--3 & $\ClMin$ least fixed point \\
\midrule
5  & Ordered ledger $\OrdLedger$ + Diamond Law (thm.) & A0$^*$, Steps~1--4 & $\WitAlg$ ordering + witness persistence \\
6  & Truth quotient $\TQ{\OrdLedger}$ & A0$^*$, Steps~1--5 & Indistinguishable $\Rightarrow$ identified \\
7  & Indistinguishability algebra ($\ObsOpen$, $\Eraser$, $\mathrm{sim}^*$) & A0$^*$, Steps~4--6 & Testability topology \\
\midrule
8  & Entropic time $\Del\Tim$ & Steps~5--7 & Append-only $\Rightarrow$ irreversible \\
9  & Entropy $S$, Budget $B$ & Steps~5, 8 & Finite resources \\
10 & Energy ledger $\EnergyLedger$, $\Efficiency$ & Steps~8--9 & Landauer cost recording \\
11 & Operational $\opnoth$ & Steps~8--10 & Budget exhaustion \\
\midrule
12 & Gauge group $\gauge$ & Steps~6, 9 & Untestable relabelings close \\
13 & Epistemic/Decision regimes & Steps~6, 8, 12 & Mode separation forced \\
14 & Exactness gate $\ExactGate$ + Split Law & Steps~5--6, 13 & Certificate-first \\
15 & Myhill--Nerode $\MNcong$ & Steps~4, 6, 12 & Futures equivalence \\
16 & Witness-path geometry ($\dsep$, $\dtime$) & Steps~4, 6, 15 & Path concat.\ $\Rightarrow$ $\triangle$-ineq.\ auto \\
17 & $\Pit$-consistent control & Steps~6, 12 & Truth blind to gauge \\
\midrule
18 & Prior-free minimax + Coupling Law & Steps~9, 14, 16 & Branch labels gauge $\Rightarrow$ minimax forced \\
19 & Self-hosting $\SelfHost = \mathcal{S}$ & Steps~3, 5, 6, 18 & System verifies itself \\
20 & Binary interface ($\SdMap$) & Steps~2, 19 & Closed-universe encoding \\
21 & Final chain $\noth \to \opnoth$ & All & End-to-end closure \\
\midrule
22 & Universal closure $\UnivOp$ & Steps~0, 4, 6, 7, 12, 16, 17 & Five-stage pipeline forced \\
23 & Closure-defect $\ClosureDefect{Q}(s)$ & Steps~6, 8, 22 & Canonical incompleteness measure \\
24 & Consciousness $\ConsProj$ & Steps~6, 18, 19, 22, 23 & Minimal persistent self-valued witness selector \\
25 & Trit field $\TritField$, valuation $\ValState$ & Steps~0, 5, 22, 23 & Component-wise terminal state \\
\midrule
26 & Continuum limit $\Omega$ & Steps~16, 9, 7 & Inverse limit of compact metric spaces \\
27 & Measure $\mu_\Pit$, $L^2(\Omega)$ & Steps~26, 12 & Sym-invariance + Kolmogorov \\
28 & Dirichlet form $\mathcal{E}$ & Steps~26, 27, 16 & Forced conductance $w = C/d^2$ \\
29 & Semigroup $\{T_t\}$, Laplacian $-\Delta$ & Steps~28, 8 & Beurling--Deny reconstruction \\
30 & Riemannian metric $g_{ij}$ & Steps~28, 26 & Sturm + Varadhan + \v{C}encov \\
31 & Complex structure $J$ & Steps~18, 30, 8 & Witness generator polar factor \\
32 & Symplectic form $\omega$ & Steps~30, 31, 5 & Replay-complete $\Rightarrow$ $d\omega = 0$ \\
33 & K\"{a}hler triple $(g, J, \omega)$ & Steps~30, 31, 32 & Single generator + $N_J = 0$ \\
\bottomrule
\end{tabular}
\end{table}

The derivation establishes that, given A0$^*$ (Completed Witnessability),
the entire reasoning apparatus---carrier, self-delimiting syntax, evaluator,
ordered witness algebra $\WitAlg$, truth quotient, observable topology,
eraser algebra, time, energy, gauge invariance, Myhill--Nerode congruence,
witness-path geometry, Bellman dynamics, self-hosting, binary interface,
universal execution law, consciousness, trit-field valuation, and the
K\"{a}hler geometric structure on the continuum limit---is \emph{forced}.
No modeling commitments survive: closed universe is forced by (W6),
commutativity is a derived sector property, prior-free minimax is forced by
absence of witnessed priors, $w = C/d^2$ is forced by scale covariance
(\cref{lem:conductance-exponent}), and $N_J = 0$ by replay-completeness
(\cref{thm:replay-nijenhuis}).  All 34 steps
(0--33) are theorems; no interpretive steps remain in the formal chain.  The
full dependency DAG is shown in \cref{fig:derivation-dag}.

% Figures for derivation (moved from main.tex)
% derivation-dag.tex — Dependency graph of the 34 derivation steps (0–33)
% Requires opoch.sty (loaded by main document)

\begin{figure}[ht]
\centering
\begin{tikzpicture}[
  % Node styles by phase (9 phases: A through I)
  origin/.style={circle, draw=opochgray, fill=opochgray!15, minimum size=8mm,
    font=\footnotesize\bfseries, inner sep=1pt},
  phaseA/.style={rounded rectangle, draw=opochblue, fill=opochblue!10,
    minimum width=14mm, minimum height=7mm, font=\footnotesize, inner sep=2pt},
  phaseB/.style={rounded rectangle, draw=opochgreen, fill=opochgreen!10,
    minimum width=14mm, minimum height=7mm, font=\footnotesize, inner sep=2pt},
  phaseC/.style={rounded rectangle, draw=opochorange, fill=opochorange!10,
    minimum width=14mm, minimum height=7mm, font=\footnotesize, inner sep=2pt},
  phaseD/.style={rounded rectangle, draw=opochred, fill=opochred!10,
    minimum width=14mm, minimum height=7mm, font=\footnotesize, inner sep=2pt},
  phaseE/.style={rounded rectangle, draw=opochblue!70!black, fill=opochblue!5,
    minimum width=14mm, minimum height=7mm, font=\footnotesize, inner sep=2pt},
  phaseF/.style={rounded rectangle, draw=opochgreen!70!black, fill=opochgreen!5,
    minimum width=14mm, minimum height=7mm, font=\footnotesize, inner sep=2pt},
  phaseG/.style={rounded rectangle, draw=opochred!70!black, fill=opochred!5,
    minimum width=14mm, minimum height=7mm, font=\footnotesize, inner sep=2pt},
  phaseH/.style={rounded rectangle, draw=opochorange!70!black, fill=opochorange!5,
    minimum width=14mm, minimum height=7mm, font=\footnotesize, inner sep=2pt},
  phaseI/.style={rounded rectangle, draw=opochpurple, fill=opochpurple!10,
    minimum width=14mm, minimum height=7mm, font=\footnotesize, inner sep=2pt},
  dep/.style={-{Stealth[length=2.5mm]}, thin, opochgray!70},
]

  % === Phase A: The Only Law (Step 0) ===
  \node[origin] (s0) at (0, 0) {$\noth$};
  \node[above=1mm of s0, font=\scriptsize\itshape, opochgray] {Step 0};

  % === Phase B: Carrier and Syntax (Steps 1–2) ===
  \node[phaseA] (s1) at (-2, -1.5) {1: $\Carrier$};
  \node[phaseA] (s2) at ( 2, -1.5) {2: $\SdMap$};

  % === Phase C: Executability (Steps 3–4) ===
  \node[phaseB] (s3) at (-2, -3.0) {3: $\UTM$};
  \node[phaseB] (s4) at ( 2, -3.0) {4: $\Del^*$};

  % === Phase D: Recording and Truth (Steps 5–7) ===
  \node[phaseC] (s5) at (-3, -4.5) {5: $\Ledger$};
  \node[phaseC] (s6) at ( 0, -4.5) {6: $\TQ{}$};
  \node[phaseC] (s7) at ( 3, -4.5) {7: $\ObsOpen$};

  % === Phase E: Time, Energy, Limits (Steps 8–11) ===
  \node[phaseD] (s8)  at (-4.5, -6.0) {8: $\Del\Tim$};
  \node[phaseD] (s9)  at (-1.5, -6.0) {9: $S, B$};
  \node[phaseD] (s10) at ( 1.5, -6.0) {10: $\EnergyLedger$};
  \node[phaseD] (s11) at ( 4.5, -6.0) {11: $\opnoth$};

  % === Phase F: Invariance and Control (Steps 12–17) ===
  \node[phaseE] (s12) at (-5, -7.5) {12: $\gauge$};
  \node[phaseE] (s13) at (-3, -7.5) {13: $\Epistemic$};
  \node[phaseE] (s14) at (-1, -7.5) {14: $\ExactGate$};
  \node[phaseE] (s15) at ( 1, -7.5) {15: $\MNcong$};
  \node[phaseE] (s16) at ( 3, -7.5) {16: $\SepDist$};
  \node[phaseE] (s17) at ( 5, -7.5) {17: $\Pit$};

  % === Phase G: Solver, Self-Hosting, Closure (Steps 18–21) ===
  \node[phaseF] (s18) at (-3, -9.0) {18: $\Bellman$};
  \node[phaseF] (s19) at ( 0, -9.0) {19: $\SelfHost$};
  \node[phaseF] (s20) at ( 3, -9.0) {20: $\SdMap$-if};

  % === Phase H: Universal Execution (Steps 22–25) ===
  \node[phaseH] (s22) at (-4.5, -10.5) {22: $\UnivOp$};
  \node[phaseH] (s23) at (-1.5, -10.5) {23: $\ClosureDefect{Q}$};
  \node[phaseH] (s24) at ( 1.5, -10.5) {24: $\ConsProj$};
  \node[phaseH] (s25) at ( 4.5, -10.5) {25: $\TritField$};

  % === Phase I: Geometric Reconstruction (Steps 26–33) ===
  \node[phaseI] (s26) at (-5.25, -12.0) {26: $\Omega$};
  \node[phaseI] (s27) at (-3.75, -12.0) {27: $L^2$};
  \node[phaseI] (s28) at (-2.25, -12.0) {28: $\mathcal{E}$};
  \node[phaseI] (s29) at (-0.75, -12.0) {29: $T_t$};
  \node[phaseI] (s30) at ( 0.75, -12.0) {30: $g$};
  \node[phaseI] (s31) at ( 2.25, -12.0) {31: $J$};
  \node[phaseI] (s32) at ( 3.75, -12.0) {32: $\omega$};
  \node[phaseI] (s33) at ( 5.25, -12.0) {33: K\"ahler};

  % === Step 21 (Closure) depends on all — bottom ===
  \node[phaseG] (s21) at (0, -13.5) {21: Closure};

  % === Dependencies ===

  % From origin (Step 0)
  \draw[dep] (s0) -- (s1);
  \draw[dep] (s0) -- (s2);

  % Phase B → Phase C
  \draw[dep] (s1) -- (s3);
  \draw[dep] (s2) -- (s3);
  \draw[dep] (s1) -- (s4);
  \draw[dep] (s3) -- (s4);

  % Phase C → Phase D
  \draw[dep] (s4) -- (s5);
  \draw[dep] (s1) -- (s5);
  \draw[dep] (s1) -- (s6);
  \draw[dep] (s5) -- (s6);
  \draw[dep] (s4) -- (s7);
  \draw[dep] (s6) -- (s7);

  % Phase D → Phase E
  \draw[dep] (s5) -- (s8);
  \draw[dep] (s7) -- (s8);
  \draw[dep] (s5) -- (s9);
  \draw[dep] (s8) -- (s9);
  \draw[dep] (s8) -- (s10);
  \draw[dep] (s9) -- (s10);
  \draw[dep] (s8) -- (s11);
  \draw[dep] (s10) -- (s11);

  % Phase E → Phase F
  \draw[dep] (s6) -- (s12);
  \draw[dep] (s9) -- (s12);
  \draw[dep] (s6) -- (s13);
  \draw[dep] (s8) -- (s13);
  \draw[dep] (s12) -- (s13);
  \draw[dep] (s5) -- (s14);
  \draw[dep] (s6) -- (s14);
  \draw[dep] (s13) -- (s14);
  \draw[dep] (s4) -- (s15);
  \draw[dep] (s6) -- (s15);
  \draw[dep] (s12) -- (s15);
  \draw[dep] (s4) -- (s16);
  \draw[dep] (s6) -- (s16);
  \draw[dep] (s15) -- (s16);
  \draw[dep] (s6) -- (s17);
  \draw[dep] (s12) -- (s17);

  % Phase F → Phase G
  \draw[dep] (s9) -- (s18);
  \draw[dep] (s14) -- (s18);
  \draw[dep] (s16) -- (s18);
  \draw[dep] (s3) -- (s19);
  \draw[dep] (s5) -- (s19);
  \draw[dep] (s6) -- (s19);
  \draw[dep] (s18) -- (s19);
  \draw[dep] (s2) -- (s20);
  \draw[dep] (s19) -- (s20);

  % Phase G → Phase H (Universal Execution)
  % 22 ← {0, 4, 6, 7, 12, 16, 17}
  \draw[dep] (s6) -- (s22);
  \draw[dep] (s7) -- (s22);
  \draw[dep] (s12) -- (s22);
  \draw[dep] (s16) -- (s22);
  \draw[dep] (s17) -- (s22);
  % 23 ← {6, 8, 22}
  \draw[dep] (s6) -- (s23);
  \draw[dep] (s8) -- (s23);
  \draw[dep] (s22) -- (s23);
  % 24 ← {6, 7, 17, 19, 22}
  \draw[dep] (s7) -- (s24);
  \draw[dep] (s17) -- (s24);
  \draw[dep] (s19) -- (s24);
  \draw[dep] (s22) -- (s24);
  % 25 ← {0, 5, 22, 23}
  \draw[dep] (s5) -- (s25);
  \draw[dep] (s22) -- (s25);
  \draw[dep] (s23) -- (s25);

  % Phase I: Geometric Reconstruction (Steps 26–33)
  % 26 ← {16, 9, 7}
  \draw[dep] (s16) -- (s26);
  \draw[dep] (s9) -- (s26);
  \draw[dep] (s7) -- (s26);
  % 27 ← {26, 12}
  \draw[dep] (s26) -- (s27);
  \draw[dep] (s12) -- (s27);
  % 28 ← {26, 27, 16}
  \draw[dep] (s26) -- (s28);
  \draw[dep] (s27) -- (s28);
  \draw[dep] (s16) -- (s28);
  % 29 ← {28, 8}
  \draw[dep] (s28) -- (s29);
  \draw[dep] (s8) -- (s29);
  % 30 ← {28, 26}
  \draw[dep] (s28) -- (s30);
  \draw[dep] (s26) -- (s30);
  % 31 ← {18, 30, 8}
  \draw[dep] (s18) -- (s31);
  \draw[dep] (s30) -- (s31);
  \draw[dep] (s8) -- (s31);
  % 32 ← {30, 31, 5}
  \draw[dep] (s30) -- (s32);
  \draw[dep] (s31) -- (s32);
  \draw[dep] (s5) -- (s32);
  % 33 ← {30, 31, 32}
  \draw[dep] (s30) -- (s33);
  \draw[dep] (s31) -- (s33);
  \draw[dep] (s32) -- (s33);

  % Step 21 (Closure) depends on all — show key edges
  \draw[dep] (s20) -- (s21);
  \draw[dep] (s25) -- (s21);
  \draw[dep] (s33) -- (s21);

  % === Legend ===
  \begin{scope}[shift={(7.5, -1.0)}]
    \node[font=\scriptsize\bfseries] at (0, 0) {Phases};
    \node[phaseA, minimum width=10mm, minimum height=5mm] at (0, -0.5) {};
    \node[font=\tiny, right] at (0.7, -0.5) {B: Carrier};
    \node[phaseB, minimum width=10mm, minimum height=5mm] at (0, -1.0) {};
    \node[font=\tiny, right] at (0.7, -1.0) {C: Executability};
    \node[phaseC, minimum width=10mm, minimum height=5mm] at (0, -1.5) {};
    \node[font=\tiny, right] at (0.7, -1.5) {D: Truth};
    \node[phaseD, minimum width=10mm, minimum height=5mm] at (0, -2.0) {};
    \node[font=\tiny, right] at (0.7, -2.0) {E: Time};
    \node[phaseE, minimum width=10mm, minimum height=5mm] at (0, -2.5) {};
    \node[font=\tiny, right] at (0.7, -2.5) {F: Invariance};
    \node[phaseF, minimum width=10mm, minimum height=5mm] at (0, -3.0) {};
    \node[font=\tiny, right] at (0.7, -3.0) {G: Solver};
    \node[phaseH, minimum width=10mm, minimum height=5mm] at (0, -3.5) {};
    \node[font=\tiny, right] at (0.7, -3.5) {H: Execution};
    \node[phaseI, minimum width=10mm, minimum height=5mm] at (0, -4.0) {};
    \node[font=\tiny, right] at (0.7, -4.0) {I: Geometry};
    \node[phaseG, minimum width=10mm, minimum height=5mm] at (0, -4.5) {};
    \node[font=\tiny, right] at (0.7, -4.5) {Closure};
  \end{scope}

\end{tikzpicture}
\caption{Dependency graph of the 34-step derivation from $\noth$ (Steps~0--33).
  Each node is a forced step; arrows show logical prerequisites.
  Colors indicate the nine derivation phases (A through I):
  The Only Law, Carrier, Executability, Truth, Time, Invariance, Solver,
  Universal Execution, and Geometry.
  Step~21 (Closure) depends on all preceding steps and appears at the bottom.}
\label{fig:derivation-dag}
\end{figure}

% fiber-shrinkage.tex — Fiber shrinkage under successive observations
% Requires opoch.sty (loaded by main document)

\begin{figure}[htbp]
\centering
\begin{tikzpicture}[
  fiberbox/.style={draw=opochblue!60, rounded corners=3pt, inner sep=4pt,
    minimum height=20mm, minimum width=16mm, fill=opochblue!3},
  dot/.style={circle, fill=#1, minimum size=2.5pt, inner sep=0pt},
  dot/.default={opochblue!70},
  dotfade/.style={circle, fill=opochgray!25, minimum size=2.5pt, inner sep=0pt},
  trans/.style={-{Stealth[length=2.5mm]}, thick, opochred!70},
  tlbl/.style={font=\scriptsize, opochred!80!black, above},
  blbl/.style={font=\scriptsize, opochorange, below},
  slbl/.style={font=\footnotesize\bfseries, opochblue!80!black, above=2mm},
]

  % ---- W0: 16 dots (4x4 grid) ----
  \node[fiberbox, minimum width=22mm, minimum height=22mm] (w0) at (0, 0) {};
  \node[slbl] at (w0.north) {$W_0$};
  \foreach \i in {0,...,3} {
    \foreach \j in {0,...,3} {
      \node[dot] at (-0.45+\i*0.3, -0.45+\j*0.3) {};
    }
  }

  % Transition 1
  \draw[trans] (1.4, 0) -- node[tlbl] {$\tau_1$ fires} node[blbl] {$\Del\Tim \geq 0$} (2.8, 0);

  % ---- W1: 8 dots (4x2 grid) ----
  \node[fiberbox, minimum width=18mm] (w1) at (4.4, 0) {};
  \node[slbl] at (w1.north) {$W_1$};
  \foreach \i in {0,...,3} {
    \foreach \j in {0,...,1} {
      \node[dot] at (3.85+\i*0.3, -0.15+\j*0.3) {};
    }
  }

  % Transition 2
  \draw[trans] (5.5, 0) -- node[tlbl] {$\tau_2$ fires} node[blbl] {$\Del\Tim \geq 0$} (6.9, 0);

  % ---- W2: 4 dots (2x2 grid) ----
  \node[fiberbox, minimum width=14mm] (w2) at (8.2, 0) {};
  \node[slbl] at (w2.north) {$W_2$};
  \foreach \i in {0,...,1} {
    \foreach \j in {0,...,1} {
      \node[dot] at (8.05+\i*0.3, -0.15+\j*0.3) {};
    }
  }

  % Transition 3
  \draw[trans] (9.2, 0) -- node[tlbl] {$\tau_3$ fires} node[blbl] {$\Del\Tim \geq 0$} (10.5, 0);

  % ---- W3: 1 dot ----
  \node[fiberbox, minimum width=10mm] (w3) at (11.5, 0) {};
  \node[slbl] at (w3.north) {$W_3$};
  \node[dot=opochgreen!80, minimum size=4pt] at (11.5, 0) {};

  % Size annotations
  \node[font=\tiny, opochgray, below=3mm of w0] {$|W_0|=16$};
  \node[font=\tiny, opochgray, below=3mm of w1] {$|W_1|=8$};
  \node[font=\tiny, opochgray, below=3mm of w2] {$|W_2|=4$};
  \node[font=\tiny, opochgray, below=3mm of w3] {$|W_3|=1$};

  % UNIQUE marker
  \node[font=\scriptsize\bfseries, opochgreen!80!black, below=7mm of w3] {$\UNIQUE$};

  % Second law annotation
  \draw[decorate, decoration={brace, amplitude=4pt, mirror}, opochorange!80]
    (-0.2, -2.0) -- (12.0, -2.0)
    node[midway, below=5pt, font=\footnotesize, opochorange!80!black]
    {Second Law: $\Del\Tim \geq 0$ at each step $\;\Longrightarrow\;$ fiber can only shrink};

\end{tikzpicture}
\caption{Fiber shrinkage under successive observations.
  Each test $\tau_i$ halves (or more) the surviving answer set $\fiber{\Ledger}$.
  The irreversible time cost $\Del\Tim \geq 0$ guarantees monotone shrinkage
  until a single answer ($\UNIQUE$) or a fundamental separator gap ($\OMEGA$) is reached.}
\label{fig:fiber-shrinkage}
\end{figure}

% diamond-law.tex — Commutative diamond: Noise and Query commute
% Requires opoch.sty (loaded by main document)

\begin{figure}[htbp]
\centering
\begin{tikzpicture}[
  node distance=25mm and 30mm,
  dnode/.style={rounded rectangle, draw=opochblue, fill=opochblue!8,
    minimum width=18mm, minimum height=9mm, font=\small, inner sep=3pt},
  darr/.style={-{Stealth[length=3mm]}, thick, opochgray!80},
  dlbl/.style={font=\small, fill=white, inner sep=2pt},
]

  % Nodes
  \node[dnode] (L) at (0, 0) {$\Ledger$};
  \node[dnode] (NL) at (-2.8, -2.5) {$\Noise(\Ledger)$};
  \node[dnode] (QL) at ( 2.8, -2.5) {$\Query(\Ledger)$};
  \node[dnode, draw=opochgreen, fill=opochgreen!8] (bot) at (0, -5.0)
    {$\Query\!\bigl(\Noise(\Ledger)\bigr) = \Noise\!\bigl(\Query(\Ledger)\bigr)$};

  % Arrows
  \draw[darr, opochred!70] (L) -- node[dlbl, left, pos=0.45]
    {$\Noise$} (NL);
  \draw[darr, opochblue!70] (L) -- node[dlbl, right, pos=0.45]
    {$\Query$} (QL);
  \draw[darr, opochblue!70] (NL) -- node[dlbl, left, pos=0.45]
    {$\Query$} (bot);
  \draw[darr, opochred!70] (QL) -- node[dlbl, right, pos=0.45]
    {$\Noise$} (bot);

  % Commutativity annotation
  \node[font=\footnotesize\itshape, opochgreen!70!black] at (0, -2.5) {commutes};

  % Forcing note
  \node[font=\scriptsize, opochgray, align=center] at (0, -6.2)
    {The diamond law ensures that encoding noise\\
     and query extraction are order-independent.};

\end{tikzpicture}
\caption{The Diamond Law: noise~$\Noise$ and query~$\Query$ commute over the
  ledger~$\Ledger$. Applying noise then querying yields the same result as
  querying then applying noise, ensuring gauge invariance of extracted truth.}
\label{fig:diamond-law}
\end{figure}

% gauge-filter.tex — Gauge filter: raw objects → G_T barrier → invariant structure
% Requires opoch.sty (loaded by main document)

\begin{figure}[htbp]
\centering
\begin{tikzpicture}[
  obj/.style={rounded rectangle, draw=#1, fill=#1!10, minimum width=10mm,
    minimum height=7mm, font=\small, inner sep=2pt},
  obj/.default={opochgray},
  blocked/.style={rounded rectangle, draw=opochred!50, fill=opochred!8,
    minimum width=10mm, minimum height=7mm, font=\small, inner sep=2pt,
    text=opochred!60},
  passed/.style={rounded rectangle, draw=opochgreen, fill=opochgreen!10,
    minimum width=10mm, minimum height=7mm, font=\small\bfseries, inner sep=2pt,
    text=opochgreen!80!black},
  barrier/.style={draw=opochorange, fill=opochorange!8, line width=1.2pt,
    minimum width=8mm, minimum height=50mm, rounded corners=2pt},
  arr/.style={-{Stealth[length=2.5mm]}, opochgray!60, thick},
  xarr/.style={-{Stealth[length=2mm]}, opochred!50, thick, dashed},
  garr/.style={-{Stealth[length=2.5mm]}, opochgreen!70, thick},
]

  % === Left side: Raw encoding-dependent objects ===
  \node[font=\footnotesize\bfseries, opochgray] at (-4.2, 2.8) {Raw objects};

  % Passing objects (encoding-invariant)
  \node[obj=opochblue] (r1) at (-4.2,  1.8) {$3$};
  \node[obj=opochblue] (r2) at (-4.2,  0.6) {$\pi$};
  \node[obj=opochblue] (r3) at (-4.2, -0.6) {$e$};

  % Blocked objects (encoding-dependent labels)
  \node[blocked] (b1) at (-4.2, -1.8) {\textsf{``three''}};
  \node[blocked] (b2) at (-4.2, -3.0) {\textsf{``drei''}};
  \node[blocked] (b3) at (-4.2, -4.2) {\textsf{``III''}};

  % === Barrier: Gauge filter G_T ===
  \node[barrier] (filt) at (0, -0.7) {};
  \node[font=\small\bfseries, opochorange!80!black, rotate=90] at (0, -0.7)
    {$\gauge$ filter};

  % Dashed vertical lines for barrier texture
  \foreach \y in {-2.8,-2.2,-1.6,-1.0,-0.4,0.2,0.8,1.4} {
    \draw[opochorange!30, thin] (-0.25, \y) -- (0.25, \y);
  }

  % === Arrows: passing through ===
  \draw[arr] (r1) -- (-0.6,  1.8);
  \draw[arr] (r2) -- (-0.6,  0.6);
  \draw[arr] (r3) -- (-0.6, -0.6);

  \draw[garr] (0.6,  1.8) -- (3.0,  1.8);
  \draw[garr] (0.6,  0.6) -- (3.0,  0.6);
  \draw[garr] (0.6, -0.6) -- (3.0, -0.6);

  % === Arrows: blocked ===
  \draw[xarr] (b1) -- (-0.6, -1.8);
  \draw[xarr] (b2) -- (-0.6, -3.0);
  \draw[xarr] (b3) -- (-0.6, -4.2);

  % X marks for blocked items
  \node[font=\large\bfseries, opochred!60] at (-0.6, -1.8) {$\times$};
  \node[font=\large\bfseries, opochred!60] at (-0.6, -3.0) {$\times$};
  \node[font=\large\bfseries, opochred!60] at (-0.6, -4.2) {$\times$};

  % === Right side: Gauge-invariant structure ===
  \node[font=\footnotesize\bfseries, opochgreen!70!black] at (4.5, 2.8)
    {Gauge-invariant};

  \node[passed] (p1) at (4.5,  1.8) {$3$};
  \node[passed] (p2) at (4.5,  0.6) {$\pi$};
  \node[passed] (p3) at (4.5, -0.6) {$e$};

  \draw[garr] (3.0,  1.8) -- (p1);
  \draw[garr] (3.0,  0.6) -- (p2);
  \draw[garr] (3.0, -0.6) -- (p3);

  % Annotation
  \node[font=\scriptsize, opochgray, align=center, text width=55mm] at (0, -5.5)
    {The gauge filter $\gauge$ blocks encoding-dependent labels\\
     and passes only structure invariant under all\\
     admissible re-encodings.};

\end{tikzpicture}
\caption{The gauge filter~$\gauge$. Raw objects enter from the left.
  Encoding-invariant quantities ($3$, $\pi$, $e$) pass through; encoding-dependent
  labels (``three'', ``drei'', ``III'') are blocked. Only gauge-invariant
  structure survives as admissible truth.}
\label{fig:gauge-filter}
\end{figure}

% observable-opens.tex — Observable opens topology on the truth quotient
% Requires opoch.sty (loaded by main document)

\begin{figure}[ht]
\centering
\begin{tikzpicture}[
  truthclass/.style={circle, fill=opochgray, minimum size=2.5pt, inner sep=0pt},
  open set label/.style={font=\scriptsize, inner sep=1pt},
]

  % --- Background: the space Pi*(L) ---
  \draw[rounded corners=12pt, draw=opochgray!50, fill=opochgray!5, line width=0.6pt]
    (-3.8, -2.8) rectangle (4.2, 3.2);
  \node[font=\footnotesize\bfseries, opochgray!80!black] at (3.4, 2.8)
    {$\TQ{\Ledger}$};

  % --- Open set O_{tau_1,a} (blue, left-center) ---
  \draw[draw=opochblue, fill=opochblue!12, line width=0.7pt, rounded corners=8pt]
    (-3.0, -0.8) .. controls (-3.2, 1.0) and (-1.8, 2.6) ..
    (-0.2, 2.4) .. controls (0.8, 2.2) and (1.0, 0.6) ..
    (0.2, -0.2) .. controls (-0.4, -1.0) and (-2.6, -1.4) ..
    cycle;
  \node[open set label, opochblue!80!black] at (-2.4, 2.2)
    {$O_{\tau_1, a}$};

  % --- Open set O_{tau_2,a} (green, right-center) ---
  \draw[draw=opochgreen, fill=opochgreen!12, line width=0.7pt, rounded corners=8pt]
    (-0.6, -1.8) .. controls (-0.8, 0.2) and (0.0, 1.8) ..
    (1.6, 2.0) .. controls (3.0, 2.0) and (3.6, 0.4) ..
    (3.2, -0.8) .. controls (2.6, -2.0) and (0.4, -2.4) ..
    cycle;
  \node[open set label, opochgreen!80!black] at (2.8, 1.8)
    {$O_{\tau_2, a}$};

  % --- Open set O_{tau_3,a} (orange, bottom) ---
  \draw[draw=opochorange, fill=opochorange!12, line width=0.7pt, rounded corners=8pt]
    (-2.4, -2.2) .. controls (-2.6, -0.8) and (-1.4, 0.4) ..
    (0.2, 0.2) .. controls (1.4, 0.0) and (2.2, -0.6) ..
    (2.0, -1.6) .. controls (1.6, -2.4) and (-1.6, -2.6) ..
    cycle;
  \node[open set label, opochorange!80!black] at (1.8, -2.2)
    {$O_{\tau_3, a}$};

  % --- Truth class points scattered across the space ---
  % In O_{tau_1,a} only
  \node[truthclass] (p1) at (-2.6, 1.4) {};
  \node[truthclass] (p2) at (-2.0, 0.6) {};
  \node[truthclass] (p3) at (-1.4, 1.8) {};

  % In O_{tau_1,a} cap O_{tau_3,a} (blue-orange intersection)
  \node[truthclass] (p4) at (-1.6, -0.4) {};
  \node[truthclass] (p5) at (-0.8, 0.0) {};

  % In O_{tau_1,a} cap O_{tau_2,a} (blue-green intersection, top)
  \node[truthclass] (p6) at (0.0, 1.6) {};
  \node[truthclass] (p7) at (0.4, 1.0) {};

  % In all three (triple intersection)
  \node[truthclass, fill=opochgray!80!black, minimum size=3pt] (p8) at (-0.1, -0.4) {};
  \node[truthclass, fill=opochgray!80!black, minimum size=3pt] (p9) at (0.4, -0.2) {};

  % In O_{tau_2,a} only
  \node[truthclass] (p10) at (2.4, 1.0) {};
  \node[truthclass] (p11) at (2.8, 0.2) {};
  \node[truthclass] (p12) at (1.8, 0.6) {};

  % In O_{tau_2,a} cap O_{tau_3,a} (green-orange intersection)
  \node[truthclass] (p13) at (1.2, -1.2) {};
  \node[truthclass] (p14) at (0.6, -1.6) {};

  % In O_{tau_3,a} only
  \node[truthclass] (p15) at (-1.0, -1.8) {};
  \node[truthclass] (p16) at (0.0, -1.4) {};

  % Outside all open sets
  \node[truthclass, fill=opochgray!40] (q1) at (3.4, -1.8) {};
  \node[truthclass, fill=opochgray!40] (q2) at (-3.2, -1.6) {};
  \node[truthclass, fill=opochgray!40] (q3) at (3.6, 2.4) {};

  % --- Annotations ---
  % Intersection label
  \draw[-{Stealth[length=2mm]}, thin, opochgray!70]
    (1.8, -2.8) node[font=\tiny, below, align=center]
      {$O_{\tau_1,a} \cap O_{\tau_2,a} \cap O_{\tau_3,a}$\\(also open)}
    -- (0.1, -0.6);

  % Union label
  \node[font=\tiny, align=center, opochgray!70!black] at (-3.4, 3.5)
    {$O_{\tau_1,a} \cup O_{\tau_2,a} \cup O_{\tau_3,a}$};
  \node[font=\tiny, opochgray!60] at (-3.4, 3.1)
    {(also open)};

  % Legend
  \begin{scope}[shift={(5.0, 1.2)}]
    \node[font=\tiny\bfseries, opochgray!80!black] at (0, 0) {Legend};
    \node[truthclass] at (-0.3, -0.5) {};
    \node[font=\tiny, right] at (0.0, -0.5) {truth class};
    \draw[draw=opochblue, fill=opochblue!12, line width=0.5pt]
      (-0.4, -1.1) rectangle (0.0, -0.8);
    \node[font=\tiny, right] at (0.0, -0.95) {$O_{\tau_1, a}$};
    \draw[draw=opochgreen, fill=opochgreen!12, line width=0.5pt]
      (-0.4, -1.5) rectangle (0.0, -1.2);
    \node[font=\tiny, right] at (0.0, -1.35) {$O_{\tau_2, a}$};
    \draw[draw=opochorange, fill=opochorange!12, line width=0.5pt]
      (-0.4, -1.9) rectangle (0.0, -1.6);
    \node[font=\tiny, right] at (0.0, -1.75) {$O_{\tau_3, a}$};
  \end{scope}

\end{tikzpicture}
\caption{Observable opens topology: open sets as testable predicates form a topology
  on the truth quotient $\TQ{\Ledger}$. Intersections and unions of observable opens
  are again open, satisfying the axioms of a topology. Individual truth classes
  (dots) are grouped by which tests separate them.}
\label{fig:observable-opens}
\end{figure}

% myhill-nerode.tex — Myhill-Nerode congruence: futures equivalence partitions states
% Requires opoch.sty (loaded by main document)

\begin{figure}[ht]
\centering
\begin{tikzpicture}[
  state/.style={circle, draw=opochgray!80, fill=white, minimum size=9mm,
    font=\footnotesize, inner sep=1pt, line width=0.6pt},
  future arrow/.style={-{Stealth[length=2mm]}, thin, #1},
  outcome/.style={rounded rectangle, draw=#1, fill=#1!15, font=\scriptsize\bfseries,
    minimum width=10mm, minimum height=6mm, inner sep=2pt},
  class label/.style={font=\scriptsize\bfseries, #1},
]

  % === Equivalence Class 1 (blue): s1, s2, s3 ===
  \begin{scope}[on background layer]
    \draw[rounded corners=10pt, draw=opochblue!50, fill=opochblue!8,
      dashed, line width=0.8pt]
      (-0.8, -0.5) rectangle (4.8, 2.5);
  \end{scope}
  \node[class label=opochblue!80!black] at (2.0, 2.2) {$[s]_1$};

  \node[state] (s1) at (0, 1.0) {$s_1$};
  \node[state] (s2) at (2, 1.0) {$s_2$};
  \node[state] (s3) at (4, 1.0) {$s_3$};

  % === Equivalence Class 2 (green): s4, s5 ===
  \begin{scope}[on background layer]
    \draw[rounded corners=10pt, draw=opochgreen!50, fill=opochgreen!8,
      dashed, line width=0.8pt]
      (-0.8, -3.5) rectangle (2.8, -1.0);
  \end{scope}
  \node[class label=opochgreen!80!black] at (1.0, -1.3) {$[s]_2$};

  \node[state] (s4) at (0, -2.2) {$s_4$};
  \node[state] (s5) at (2, -2.2) {$s_5$};

  % === Equivalence Class 3 (orange): s6, s7, s8 ===
  \begin{scope}[on background layer]
    \draw[rounded corners=10pt, draw=opochorange!50, fill=opochorange!8,
      dashed, line width=0.8pt]
      (3.2, -3.5) rectangle (8.8, -1.0);
  \end{scope}
  \node[class label=opochorange!80!black] at (6.0, -1.3) {$[s]_3$};

  \node[state] (s6) at (4, -2.2) {$s_6$};
  \node[state] (s7) at (6, -2.2) {$s_7$};
  \node[state] (s8) at (8, -2.2) {$s_8$};

  % === Terminal outcomes (right side) ===
  \node[outcome=opochgreen] (unique) at (10.5, 1.0) {$\UNIQUE$};
  \node[outcome=opochorange] (omega)  at (10.5, -2.2) {$\OMEGA$};

  % === Futures arrows from Class 1 states -> UNIQUE ===
  \draw[future arrow=opochblue!70] (s1) -- ++(1.8, 0)
    node[font=\tiny, above, pos=0.4, opochblue!70!black] {$\sigma_1$};
  \draw[future arrow=opochblue!70] (s2) -- ++(1.8, 0)
    node[font=\tiny, above, pos=0.4, opochblue!70!black] {$\sigma_1$};
  \draw[future arrow=opochblue!70] (s3) -- ++(1.8, 0)
    node[font=\tiny, above, pos=0.4, opochblue!70!black] {$\sigma_1$};

  % Merge arrows toward UNIQUE
  \draw[future arrow=opochblue!70] (s1.east) ++(1.8, 0) -- (unique.west);
  \draw[future arrow=opochblue!70] (s2.east) ++(1.8, 0) -- (unique.west);
  \draw[future arrow=opochblue!70] (s3.east) ++(1.8, 0) -- (unique.west);

  % === Futures arrows from Class 2 states -> UNIQUE ===
  \draw[future arrow=opochgreen!70] (s4) -- ++(2.0, 0)
    node[font=\tiny, above, pos=0.4, opochgreen!70!black] {$\sigma_2$};
  \draw[future arrow=opochgreen!70] (s5) -- ++(2.0, 0)
    node[font=\tiny, above, pos=0.4, opochgreen!70!black] {$\sigma_2$};

  % Bend class 2 arrows up to UNIQUE
  \draw[future arrow=opochgreen!70] (s4.east) ++(2.0, 0) -- (unique.south west);
  \draw[future arrow=opochgreen!70] (s5.east) ++(2.0, 0) -- (unique.south west);

  % === Futures arrows from Class 3 states -> OMEGA ===
  \draw[future arrow=opochorange!70] (s6) -- ++(0.5, 0)
    node[font=\tiny, above, pos=0.5, opochorange!70!black] {$\sigma_1$};
  \draw[future arrow=opochorange!70] (s7) -- ++(0.5, 0)
    node[font=\tiny, above, pos=0.5, opochorange!70!black] {$\sigma_1$};
  \draw[future arrow=opochorange!70] (s8) -- ++(0.5, 0)
    node[font=\tiny, above, pos=0.5, opochorange!70!black] {$\sigma_1$};

  \draw[future arrow=opochorange!70] (s6.east) ++(0.5, 0) -- (omega.west);
  \draw[future arrow=opochorange!70] (s7.east) ++(0.5, 0) -- (omega.west);
  \draw[future arrow=opochorange!70] (s8.east) ++(0.5, 0) -- (omega.west);

  % === Annotation: equivalence meaning ===
  \node[font=\tiny, align=center, opochgray!70!black] at (4.0, -4.3)
    {States in the same class share identical futures:\\
     $s \MNcong s' \;\Leftrightarrow\; \forall \sigma:\;
     \mathrm{outcome}(s, \sigma) = \mathrm{outcome}(s', \sigma)$};

  % === Bidirectional equivalence arrows within classes ===
  \draw[{Stealth[length=1.5mm]}-{Stealth[length=1.5mm]}, densely dotted,
    opochblue!50, thin] (s1) -- (s2);
  \draw[{Stealth[length=1.5mm]}-{Stealth[length=1.5mm]}, densely dotted,
    opochblue!50, thin] (s2) -- (s3);

  \draw[{Stealth[length=1.5mm]}-{Stealth[length=1.5mm]}, densely dotted,
    opochgreen!50, thin] (s4) -- (s5);

  \draw[{Stealth[length=1.5mm]}-{Stealth[length=1.5mm]}, densely dotted,
    opochorange!50, thin] (s6) -- (s7);
  \draw[{Stealth[length=1.5mm]}-{Stealth[length=1.5mm]}, densely dotted,
    opochorange!50, thin] (s7) -- (s8);

\end{tikzpicture}
\caption{Myhill--Nerode congruence: states are partitioned into equivalence classes
  by futures equivalence $\MNcong$. States within the same class (dashed regions)
  produce identical terminal outcomes ($\UNIQUE$ or $\OMEGA$) under every test
  sequence~$\sigma$. The quotient automaton has exactly one state per class.}
\label{fig:myhill-nerode}
\end{figure}

% separator-geometry.tex — Separator geometry: metric on the truth quotient
% Requires opoch.sty (loaded by main document)

\begin{figure}[ht]
\centering
\begin{tikzpicture}[
  truthpt/.style={circle, fill=#1, draw=black!30, minimum size=5pt,
    inner sep=0pt, line width=0.4pt},
  truthpt/.default={opochblue!70},
  edge label/.style={font=\tiny, fill=white, inner sep=1pt, rounded corners=1pt},
  region label/.style={font=\scriptsize\itshape, #1},
]

  % === Curved surface (manifold sketch) ===
  % Draw a subtle grid to suggest a curved surface
  \begin{scope}[opochblue!15, line width=0.3pt]
    % Horizontal curves (slight curvature)
    \draw (-4.0, -1.5) .. controls (-2.0, -1.3) and (0.0, -1.8) ..
      (2.0, -1.2) .. controls (3.5, -0.8) and (4.5, -1.0) .. (5.5, -0.6);
    \draw (-4.0, -0.5) .. controls (-2.0, -0.2) and (0.0, -0.8) ..
      (2.0, -0.1) .. controls (3.5, 0.3) and (4.5, 0.0) .. (5.5, 0.4);
    \draw (-4.0, 0.5) .. controls (-2.0, 0.8) and (0.0, 0.2) ..
      (2.0, 0.9) .. controls (3.5, 1.3) and (4.5, 1.0) .. (5.5, 1.4);
    \draw (-4.0, 1.5) .. controls (-2.0, 1.9) and (0.0, 1.2) ..
      (2.0, 1.9) .. controls (3.5, 2.3) and (4.5, 2.0) .. (5.5, 2.4);
    % Vertical curves
    \draw (-3.5, -1.8) .. controls (-3.4, -0.5) and (-3.6, 0.5) .. (-3.5, 1.8);
    \draw (-1.5, -1.6) .. controls (-1.4, -0.3) and (-1.6, 0.7) .. (-1.5, 1.9);
    \draw (0.5, -1.9) .. controls (0.6, -0.6) and (0.4, 0.4) .. (0.5, 1.5);
    \draw (2.5, -1.2) .. controls (2.6, 0.0) and (2.4, 1.0) .. (2.5, 2.0);
    \draw (4.5, -0.9) .. controls (4.6, 0.2) and (4.4, 1.2) .. (4.5, 2.2);
  \end{scope}

  % === High curvature region (left, closely spaced points) ===
  \begin{scope}[on background layer]
    \draw[rounded corners=6pt, draw=opochred!30, fill=opochred!6,
      densely dashed, line width=0.5pt]
      (-3.2, -0.8) rectangle (-0.8, 1.2);
  \end{scope}

  \node[truthpt=opochred!60] (p1) at (-2.8, 0.8) {};
  \node[truthpt=opochred!60] (p2) at (-2.2, 0.2) {};
  \node[truthpt=opochred!60] (p3) at (-1.4, 0.6) {};
  \node[truthpt=opochred!60] (p4) at (-2.0, -0.4) {};

  \node[font=\tiny, above=1mm] at (p1) {$[p_1]$};
  \node[font=\tiny, below left=0.5mm] at (p2) {$[p_2]$};
  \node[font=\tiny, above right=0.5mm] at (p3) {$[p_3]$};
  \node[font=\tiny, below=1mm] at (p4) {$[p_4]$};

  % Edges with distances (high curvature = small metric distances but high cost)
  \draw[opochred!60, line width=0.5pt] (p1) -- (p2)
    node[edge label, pos=0.5] {$\SepDist\!=\!3$};
  \draw[opochred!60, line width=0.5pt] (p2) -- (p3)
    node[edge label, pos=0.5] {$\SepDist\!=\!4$};
  \draw[opochred!60, line width=0.5pt] (p2) -- (p4)
    node[edge label, pos=0.5] {$\SepDist\!=\!2$};
  \draw[opochred!60, line width=0.5pt] (p3) -- (p4)
    node[edge label, pos=0.4] {$\SepDist\!=\!5$};

  % Curvature annotation
  \node[region label=opochred!80!black, align=center] at (-2.0, -1.4)
    {High curvature $\SepCurv \gg 0$\\(expensive separators)};

  % === Low curvature region (right, spread out points) ===
  \begin{scope}[on background layer]
    \draw[rounded corners=6pt, draw=opochblue!30, fill=opochblue!6,
      densely dashed, line width=0.5pt]
      (1.0, -1.0) rectangle (5.2, 2.2);
  \end{scope}

  \node[truthpt=opochblue!70] (q1) at (1.6, 1.6) {};
  \node[truthpt=opochblue!70] (q2) at (3.8, 1.8) {};
  \node[truthpt=opochblue!70] (q3) at (4.6, 0.0) {};
  \node[truthpt=opochblue!70] (q4) at (2.4, -0.4) {};

  \node[font=\tiny, above=1mm] at (q1) {$[q_1]$};
  \node[font=\tiny, above=1mm] at (q2) {$[q_2]$};
  \node[font=\tiny, right=1mm] at (q3) {$[q_3]$};
  \node[font=\tiny, below=1mm] at (q4) {$[q_4]$};

  % Edges with distances (low curvature = cheap separators)
  \draw[opochblue!60, line width=0.5pt] (q1) -- (q2)
    node[edge label, pos=0.5] {$\SepDist\!=\!1$};
  \draw[opochblue!60, line width=0.5pt] (q2) -- (q3)
    node[edge label, pos=0.5] {$\SepDist\!=\!1$};
  \draw[opochblue!60, line width=0.5pt] (q3) -- (q4)
    node[edge label, pos=0.5] {$\SepDist\!=\!1$};
  \draw[opochblue!60, line width=0.5pt] (q1) -- (q4)
    node[edge label, pos=0.5] {$\SepDist\!=\!1$};

  % Curvature annotation
  \node[region label=opochblue!80!black, align=center] at (3.4, -1.4)
    {Low curvature $\SepCurv \approx 0$\\(cheap separators)};

  % === Cross-region edge ===
  \draw[opochgray!50, line width=0.5pt, densely dotted] (p3) -- (q1)
    node[edge label, pos=0.5] {$\SepDist\!=\!7$};

  % === Axis label ===
  \node[font=\footnotesize\bfseries, opochgray!80!black] at (1.0, 2.8)
    {$(\TQ{\Ledger},\; \SepDist)$};

\end{tikzpicture}
\caption{Separator geometry: the cost of distinguishing truth classes induces
  a metric~$\SepDist$ on the quotient space $\TQ{\Ledger}$. Regions where
  separators are expensive exhibit high curvature~$\SepCurv$, while regions
  with cheap single-test separators are metrically flat.}
\label{fig:separator-geometry}
\end{figure}

% self-hosting.tex — Self-hosting: the system verifies its own derivation
% Requires opoch.sty (loaded by main document)

\begin{figure}[ht]
\centering
\begin{tikzpicture}[
  box/.style={rectangle, draw=#1, fill=#1!10, rounded corners=4pt,
    minimum width=32mm, minimum height=10mm, font=\small,
    line width=0.7pt, align=center, inner sep=4pt},
  box/.default={opochgray},
  loop arrow/.style={-{Stealth[length=2.5mm]}, line width=1pt, #1},
  step label/.style={font=\tiny\itshape, align=center, text width=25mm, #1},
]

  % === System S at the top ===
  \node[box=opochblue] (kernel) at (0, 0)
    {System $\mathcal{S}$};
  \node[step label=opochblue!70!black, above=2mm of kernel]
    {\Opoch{} system};

  % === Step 1: Encode derivation ===
  \node[box=opochgray] (encode) at (0, -2.5)
    {Encode derivation as\\$P_0, P_1, \ldots, P_{33}$};

  % === Step 2: Run S on each P_i ===
  \node[box=opochorange] (run) at (0, -5.0)
    {Run $\mathcal{S}$ on each $P_i$};

  % Individual step indicators
  \begin{scope}[shift={(0, -5.0)}]
    \foreach \i in {0,...,5} {
      \node[circle, draw=opochorange!50, fill=opochorange!15,
        minimum size=4mm, font=\tiny, inner sep=0pt]
        at ({-2.5 + \i*1.0}, -1.2) {$P_{\i}$};
    }
    \node[font=\tiny, opochorange!60] at (4.0, -1.2) {$\cdots\; P_{33}$};
  \end{scope}

  % === Step 3: All produce UNIQUE ===
  \node[box=opochgreen] (verify) at (0, -8.5)
    {All produce $\UNIQUE$};

  % Checkmarks for visual confirmation
  \begin{scope}[shift={(0, -8.5)}]
    \foreach \x in {-1.0, 0, 1.0} {
      \node[font=\small, opochgreen!70!black] at (\x, -0.9) {\checkmark};
    }
  \end{scope}

  % === Downward arrows ===
  \draw[loop arrow=opochblue!70] (kernel) -- (encode)
    node[step label=opochgray, pos=0.5, right=3mm]
      {Step 1: encode each\\derivation step\\as a proposition};

  \draw[loop arrow=opochorange!70] (encode) -- (run)
    node[step label=opochgray, pos=0.5, right=3mm]
      {Step 2: submit\\to system};

  \draw[loop arrow=opochgreen!70] (run) -- ++(0, -1.2)
    -- (verify.north -| 0, 0)
    node[step label=opochgray, pos=0.3, right=3mm]
      {Step 3: each $P_i$\\resolves to $+1$};

  % === Feedback loop: back to System ===
  \draw[loop arrow=opochgreen!60!black, line width=1.2pt]
    (verify.west) -- ++(-3.2, 0) -- ++(0, 8.5)
    -- (kernel.west)
    node[step label=opochgreen!70!black, pos=0.5, left=3mm, align=center]
      {$\FixPt(\mathcal{S}) = \mathcal{S}$\\[2pt]
       \footnotesize Self-hosting\\
       \footnotesize fixed point};

  % === Central emphasis ===
  \node[font=\tiny, opochgray!60, align=center] at (3.8, -4.2)
    {The system treats its own\\derivation steps as input\\and verifies each one.};

  % === Receipt box ===
  \node[rectangle, draw=opochgreen!40, fill=opochgreen!5,
    rounded corners=2pt, font=\tiny, inner sep=3pt, align=left]
    at (3.8, -9.5)
    {Receipt:\\
     $V(P_i, w_i) = +1\;\;\forall\, i \in \{0,\ldots,33\}$\\
     $\Rightarrow\;\FixPt(\mathcal{S}) = \mathcal{S}$};

\end{tikzpicture}
\caption{Self-hosting: the system~$\mathcal{S}$ verifies its own derivation
  via~$\FixPt(\mathcal{S}) = \mathcal{S}$. Each derivation step $P_0, \ldots, P_{33}$
  is encoded as a proposition and fed to~$\mathcal{S}$; all 34 produce $\UNIQUE$,
  closing the self-referential loop.}
\label{fig:self-hosting}
\end{figure}

